\documentclass[a4paper, oneside, 13pt]{report}
\usepackage{vntex}
% \usepackage[english,vietnam]{babel}
\usepackage[utf8]{inputenc}

% \usepackage[utf8]{inputenc}
% %\usepackage[francais]{babel}
\usepackage{a4wide,amssymb,epsfig,latexsym,array,hhline,fancyhdr}
\usepackage[normalem]{ulem}
\usepackage{soul}
\usepackage{tabto}
\usepackage[makeroom]{cancel}
\usepackage{amsmath}
\usepackage{amsthm}
\usepackage{multicol,longtable,amscd}
\usepackage{diagbox}%Make diagonal lines in tables
\usepackage{booktabs}
\usepackage{alltt}
\usepackage[framemethod=tikz]{mdframed}% For highlighting paragraph backgrounds
\usepackage{caption,subcaption}

\usepackage{lastpage}
\usepackage[lined,boxed,commentsnumbered]{algorithm2e}
\usepackage{enumerate}
\usepackage{color}
\usepackage{graphicx}							% Standard graphics package
\usepackage{array}
\usepackage{tabularx, caption}
\usepackage{multirow}
\usepackage{multicol}
\usepackage{rotating}
\usepackage{graphics}
\usepackage{geometry}
\usepackage{setspace}
\usepackage{epsfig}
\usepackage{tikz}
\usepackage{times}
\usetikzlibrary{arrows,snakes,backgrounds}
\usepackage[unicode]{hyperref}
\hypersetup{urlcolor=blue,linkcolor=black,citecolor=black,colorlinks=true} 
%\usepackage{pstcol} 								% PSTricks with the standard color package

\usepackage[normalem]{ulem}


\newtheorem{theorem}{{\bf Định lý}}
\newtheorem{property}{{\bf Tính chất}}
\newtheorem{proposition}{{\bf Mệnh đề}}
\newtheorem{corollary}[proposition]{{\bf Hệ quả}}
\newtheorem{lemma}[proposition]{{\bf Bổ đề}}
\theoremstyle{definition}
\newtheorem{exer}{Bài toán}

\def\thesislayout{	% A4: 210 × 297
	\geometry{
		a4paper,
		total={160mm,240mm},  % fix over page
		left=3cm,
		top=2cm,
		right=2cm,
		bottom=2cm, 
	}
}
\thesislayout

%\usepackage{fancyhdr}
\setlength{\headheight}{40pt}
\pagestyle{fancy}
\fancyhead{} % clear all header fields
\fancyhead[L]{
 \begin{tabular}{rl}
    \begin{picture}(25,15)(0,0)
    \put(0,-8){\includegraphics[width=8mm, height=8mm]{img/hcmut.png}}
    %\put(0,-8){\epsfig{width=10mm,figure=hcmut.eps}}
   \end{picture}&
	% \includegraphics[width=8mm, height=8mm]{hcmut.png} & %
	\begin{tabular}{l}
		\textbf{\bf \ttfamily Trường Đại Học Bách Khoa Tp.Hồ Chí Minh}\\
		\textbf{\bf \ttfamily Khoa Khoa Học \& Kỹ Thuật Máy Tính}
	\end{tabular} 	
 \end{tabular}
}
\fancyhead[R]{
	\begin{tabular}{l}
		\tiny \bf \\
		\tiny \bf 
	\end{tabular}  }
\fancyfoot{} % clear all footer fields
\fancyfoot[L]{\scriptsize \ttfamily Đồ án chuyên ngành}
\fancyfoot[R]{\scriptsize \ttfamily Trang {\thepage}/\pageref{LastPage}}
\renewcommand{\headrulewidth}{0.3pt}
\renewcommand{\footrulewidth}{0.3pt}

% %%%
\setcounter{secnumdepth}{4}
\setcounter{tocdepth}{3}
\makeatletter
\newcounter {subsubsubsection}[subsubsection]
\renewcommand\thesubsubsubsection{\thesubsubsection .\@alph\c@subsubsubsection}
\newcommand\subsubsubsection{\@startsection{subsubsubsection}{4}{\z@}%
                                     {-3.25ex\@plus -1ex \@minus -.2ex}%
                                     {1.5ex \@plus .2ex}%
                                     {\normalfont\normalsize\bfseries}}
\newcommand*\l@subsubsubsection{\@dottedtocline{3}{10.0em}{4.1em}}
\newcommand*{\subsubsubsectionmark}[1]{}
\makeatother

\everymath{\color{black}}%make in-line maths symbols blue to read/check easily

\sloppy
\captionsetup[figure]{labelfont={small,bf},textfont={small,it},belowskip=-1pt,aboveskip=-9pt}
%space remove between caption, figure, and text
\captionsetup[table]{labelfont={small,bf},textfont={small,it},belowskip=-1pt,aboveskip=7pt}

\setlength{\floatsep}{5pt plus 2pt minus 2pt}
\setlength{\textfloatsep}{5pt plus 2pt minus 2pt}
\setlength{\intextsep}{10pt plus 2pt minus 2pt}

\thesislayout

\newcommand{\setupfont}[1]{\fontsize{#1}{#1}\selectfont}

\DeclareMathOperator*{\argmax}{arg\,max}
\DeclareMathOperator*{\argmin}{arg\,min}

\newcommand{\norm}[1]{\lVert #1 \rVert}
\newcommand{\bignorm}[1]{\bigg \lVert #1 \bigg \rVert}

\newcommand{\recomplex}[1]{#1}
\newcommand{\imcomplex}[1]{#1i}
\newcommand{\complex}[2]{#1 + i#2}
\newcommand{\polarcomplex}[2]{#1e^{#2i}}
\newcommand{\tricomplex}[2]{\cos(#1) + i\sin(#2)}

\newcommand{\cartesiancoord}{Oft}
\newcommand{\complexcoord}{O \Re \Im}

\newcommand{\ifc}{\text{nếu }}
\newcommand{\otherwise}{\text{các trường hợp khác}}

\newcommand{\scalemath}[2]{\scalebox{#1}{\mbox{\ensuremath{\displaystyle #2}}}}

\newcommand{\definition}[1]{\textbf{\textit{#1}}}

\newcommand{\formularef}[1]{(\ref{#1})}
\newcommand{\enumentry}[1]{\textit{#1}}
\newcommand{\name}[1]{\text{#1}}

\newcommand{\figref}[1]{Hình~\ref{#1}}
\newcommand{\tableref}[1]{Bảng~\ref{#1}}
\newcommand{\labelref}[2]{#1~\ref{#2}}


\newcommand{\chapterref}[1]{Chương~\ref{#1}}
\newcommand{\sectionref}[1]{Mục~\ref{#1}}
\newcommand{\subsectionref}[1]{Tiểu mục~\ref{#1}}
\newcommand{\appendixref}[1]{Phụ lục~\ref{#1}}

\newcommand{\urlsrc}[1]{(nguồn~\url{#1})}
\newcommand{\citesrc}[1]{(nguồn~\cite{#1})}

\newcommand{\srccaption}[2]{\caption[#1]{#1 #2}}

\newcommand{\numberedeq}{\addtocounter{equation}{1}\tag{\theequation}}
\newcommand{\bigO}{\name{O}}

\newcommand{\Waveform}{\text{Waveform}}
\newcommand{\Spectrum}{\text{Spectrum}}
\newcommand{\Spectrogram}{\text{Spectrogram}}

\newcommand{\waveform}{\text{waveform}}
\newcommand{\spectrum}{\text{spectrum}}
\newcommand{\spectrogram}{\text{spectrogram}}

\newcommand{\subsub}[1]{\noindent \textbf{#1}.}

% \newcommand{\theorem}{Theorem}

\usepackage{fancyhdr}
\renewcommand{\sectionmark}[1]{\markright{\thesection\ #1}}
\fancypagestyle{fancy}{
	\fancyhf{}
	\renewcommand{\headrulewidth}{0.4pt}
	\renewcommand{\footrulewidth}{0pt}
	\fancyhead[OR]{\bfseries \nouppercase{\rightmark}}
	\fancyhead[EL]{\bfseries \nouppercase{\leftmark}}
	\fancyfoot[OR, EL]{\thepage}
}

\fancypagestyle{plain}{
	\fancyhf{}
	\renewcommand{\headrulewidth}{0pt}
	\renewcommand{\footrulewidth}{0pt}
	\fancyfoot[OR, EL]{\thepage}
}
\fancypagestyle{addpagenumbersforpdfimports}{
	\fancyhead{}
	\renewcommand{\headrulewidth}{0pt}
	\fancyfoot{}
	\fancyfoot[R]{\thepage}
}


\usepackage{tikz}
\usetikzlibrary{calc}
\newcommand\HRule{\rule{\textwidth}{1pt}}

\begin{document}
	
	\tolerance=1
	\emergencystretch=\maxdimen
	\hyphenpenalty=10000
	\hbadness=10000
	
	\setlength{\abovedisplayskip}{-2.5mm}
	\setlength{\belowdisplayskip}{0pt}
	\setlength{\abovedisplayshortskip}{-2.5mm}
	%\setlength{\belowdisplayshortskip}{-0.5pt}

\begin{titlepage}
	\begin{tikzpicture}[remember picture, overlay]
		\draw[line width = 4pt] ($(current page.north west) + (1in + 5,-1in+30)$) rectangle ($(current page.south east) + (-1in+30,1in-20)$);
		\draw[line width = 1pt] ($(current page.north west) + (1in,-1in+35)$) rectangle ($(current page.south east) + (-1in+35,1in-25)$);
	  \end{tikzpicture}

	
	\begin{center}
		\setupfont{14pt}
	ĐẠI HỌC QUỐC GIA THÀNH PHỐ HỒ CHÍ MINH \\
	TRƯỜNG ĐẠI HỌC BÁCH KHOA \\
	KHOA KHOA HỌC VÀ KỸ THUẬT MÁY TÍNH 
	\end{center}

\vspace{1cm}

\begin{figure}[h!]
\begin{center}
\includegraphics[width=50mm]{img/hcmut.png}
\end{center}
\end{figure}


\begin{center}
\setupfont{15pt}
	\begin{tabular}{c}
  %       \centering{\textbf{{\Large BÁO CÁO }}}
  %       \newline
		% \multicolumn{1}{l}{\textbf{{\Large ĐỒ ÁN CHUYÊN NGÀNH }}}
        \textbf{{BÁO CÁO}} \\
        \textbf{{ĐỒ ÁN CHUYÊN NGÀNH}}
	\end{tabular}
\end{center}

\vspace*{0.5cm}

\begin{center}
\setupfont{17pt}
	\begin{tabular}{c}
	\textbf{{XÂY DỰNG HỆ THỐNG QUẢN LÝ}}\\
	\textbf{{CỬA HÀNG THỜI TRANG DỰA TRÊN}}\\
	\textbf{{VIỆC TỰ ĐỘNG HÓA QUY TRÌNH NGHIỆP VỤ}}
	\end{tabular}
\end{center}

\vspace{0.5cm}


\begin{center}
	\setupfont{14pt}
	\begin{tabular}{rl}
		\textbf{Ngành:} & Khoa học máy tính
	\end{tabular}
\end{center}

\vspace{1cm}

\begin{table}[h]
	\setupfont{13pt}
\begin{tabular}{rll}
\hspace{5 cm}	& \textbf{HỘI ĐỒNG:} 	& Công nghệ phần mềm  \\
				& \textbf{GVHD:} 		& ThS. Trương Thị Thái Minh\\
				&						& TS. Nguyễn An Khương\\
				& \textbf{TKHĐ:}		& Trương Thị Thái Minh\\
				&&\multicolumn{1}{c}{\bf---o0o---}\vspace{2mm}\\
				& \textbf{SVTH 1:} 		& Bùi Lương Vinh Hiển (1913380)       \\
				& \textbf{SVTH 2:} 		& Trần Tuấn Phong (1914641)       \\
				& \textbf{SVTH 3:} 		& Trần Nguyễn Hữu Thọ (1915347)       \\
\end{tabular}
\end{table}

\begin{center}
{\setupfont{12pt} Thành phố Hồ Chí Minh, 12/2022}
\end{center}
\end{titlepage}


\newpage
\newpage
\setlength{\parskip}{2.5mm}
\thispagestyle{plain}
	\begin{center}
		\section*{Lời cam đoan}
		\addcontentsline{toc}{chapter}{Lời cam đoan}
	\end{center}
	
	Tôi xin cam đoan đây là công trình nghiên cứu của riêng tôi dưới sự hướng dẫn của TS. Nguyễn An Khương, TS. Nguyễn Tiến Thịnh, KS. Nguyễn Tấn Đức, KS. Nguyễn Hữu Hồng Huy. Nội dung nghiên cứu và các kết quả đều là trung thực và chưa từng được công bố trước đây. Các số liệu được sử dụng cho quá trình phân tích, nhận xét được chính tôi thu thập từ nhiều nguồn khác nhau và sẽ được ghi rõ trong phần tài liệu tham khảo. Ngoài ra, tôi cũng có sử dụng một số nhận xét, đánh giá và số liệu của các tác giả khác, cơ quan tổ chức khác. Tất cả đều có trích dẫn và chú thích nguồn gốc. Nếu phát hiện có bất kì sự gian lận nào, tôi xin hoàn toàn chịu trách nhiệm về nội dung luận văn tốt nghiệp của mình. Trường đại học Bách Khoa thành phố Hồ Chí Minh không liên quan đến những vi phạm tác quyền, bản quyền do tôi gây ra trong quá trình thực hiện.
	
	\vspace{1cm}
	
	\begin{center}
		\begin{tabular}{c c c}
			KÝ TÊN: \rule{4cm}{0.15mm} &\hspace{4cm}& NGÀY: \rule{4cm}{0.15mm}
		\end{tabular}
	\end{center}
\setlength{\parskip}{0mm}

\newpage


\tableofcontents

%%%%%%%%%%%%%%%%%%%%%%%%%%%%%%%%%%%%

% Giới thiệu
\chapter{Tổng quan}
% !TeX root = ..\main.tex

%%%%%%%%%%%%%%%%%%%%%%%%%
\section{Giới thiệu đề tài}

\hspace{0.5cm}Quy trình nghiệp vụ là một chuỗi các hành động diễn ra có thứ tự để cùng nhau hoàn thành một mục tiêu hoặc đạt được mục đích đề ra. Quy trình nghiệp vụ là một thành phần quan trọng trong quá trình hoạt động của một doanh nghiệp, chúng đại diện cho một trong những tài sản cốt lõi của các tổ chức. Chúng có tác động trực tiếp đến sức hấp dẫn của sản phẩm và dịch vụ, ảnh hưởng đến trải nghiệm của khách hàng và cuối cùng là ảnh hưởng đến doanh thu của doanh nghiệp. Các quy trình nghiệp vụ phối hợp cùng nhau để đáp ứng các nhu cầu từ người dùng, do đó các quy trình này cần được thiết kế đúng với nhu cầu của người dùng và hệ thống. Cụ thể, quy trình nghiệp vụ xác định các công việc và trách nhiệm của mỗi thành phần trong doanh nghiệp, thứ tự thực thi của các thành phần và hướng xử lý khi có lỗi. Các quy trình nghiệp vụ là hệ thống quan trọng trong tổ chức và trong mạng lưới cung ứng liên tổ chức. Do đó, bất kỳ lỗi quy trình nào cũng có thể khiến hoạt động của công ty và toàn bộ hệ sinh thái bị đình trệ.\\

Quy trình nghiệp vụ có tầm quan trọng là vậy, tuy nhiên quản lý việc vận hành quy trình nghiệp vụ là một bài toán không hề dễ dàng. Với việc các doanh nghiệp ngày càng mở rộng quy mô kinh doanh, cùng sự tăng lên của tập người dùng, các quy trình nghiệp vụ cũng dần lớn lên, và thường phải thay đổi theo sự thay đổi của người dùng và thị trường. Sự thay đổi nhanh của các quy trình nghiệp vụ đòi hỏi một phương pháp quản lý hiệu quả hơn cho các quy trình này. Xuất phát từ nhu cầu đó, một bộ công cụ, kỹ thuật và phương pháp được sinh ra để hỗ trợ quản lý vòng đời quy trình nghiệp vụ đã xuất hiện và được sử dụng trong hai thập kỷ qua. Nó được gọi là \textbf{Quản lý quy trình nghiệp vụ (BPM)}. BPM được tin dùng cho nhiều lĩnh vực khác nhau, bao gồm kỹ thuật công nghiệp, quản lý hoạt động, quản lý chất lượng, quản lý nguồn nhân lực, quản trị doanh nghiệp, khoa học máy tính và kỹ thuật, hệ thống thông tin.\\

Trong quá trình thực thi quy trình nghiệp vụ, những quy trình có thể cần được cải thiện và thay đổi. Khi sửa đổi quy trình, doanh nghiệp cần đảm bảo chương trình chạy đúng như quy trình mới. Việc thay đổi mã nguồn bằng tay để thay đổi luồng thực thi chương trình sẽ tiêu tốn rất nhiều thời gian và công sức; tuy vậy, sai lầm lại rất dễ xảy ra. Từ đây, nhu cầu áp dụng công nghệ vào việc tự động hóa các quy trình được đặt ra. Để đáp ứng nhu cầu này, phương pháp \textbf{Tự động hóa quy trình nghiệp vụ (BPA)} được ra đời. Thay vì dùng con người, phương pháp này sử dụng công nghệ và phần mềm để tự động hóa các quy trình nghiệp vụ. Trái ngược với các loại tự động hóa khác, các giải pháp BPA có xu hướng phức tạp, được kết nối với nhiều hệ thống công nghệ thông tin  và được điều chỉnh cụ thể theo nhu cầu của một tổ chức. Các tổ chức thường áp dụng BPA như một phần của chiến lược chuyển đổi kỹ thuật số, nhằm hợp lý hóa quy trình nghiệp vụ của mình và đảm bảo hiệu quả hoạt động. Tự động hóa quy trình nghiệp vụ có thể giải phóng thời gian và nguồn lực cho doanh nghiệp rất nhiều.\\

Với các yêu cầu về thiết kế và quản lý quy trình nghiệp vụ như trên, Lược đồ mô hình hóa quy trình nghiệp vụ (BPMN) là một lựa chọn rất hữu ích. Đây là ngôn ngữ hướng đồ thị, cung cấp cho doanh nghiệp khả năng thiết kế các quy trình nghiệp vụ của mình dưới dạng lược đồ, đồng thời đặt ra tiêu chuẩn chung cho các doanh nghiệp truyền đạt các quy trình nghiệp vụ. Việc biểu diễn các quy trình nghiệp vụ theo dạng lược đồ giúp cho doanh nghiệp quản lý tốt hơn hoạt động kinh doanh của mình khi mọi phòng ban đều có thể hiểu ngôn ngữ ở dạng lược đồ. Để thực thi quy trình nghiệp vụ được thiết kể bởi BPMN, chúng ta cần sử dụng đến Ngôn ngữ thực thi quy trình nghiệp vụ (BPEL). BPMN đóng vai trò phác thảo các quy trình nghiệp vụ, là công cụ để các thành phần trong doanh nghiệp có thể trao đổi và thiết kế quy trình nghiệp vụ. Sau khi thiết kế BPMN, BPEL là ngôn ngữ để vận hành các quy trình nghiệp vụ đó. Với BPMN và BPEL, ta có thể dễ dàng tự động hóa quy trình nghiệp vụ, bảo đảm quy trình nghiệp vụ được vận hành chính xác, và việc cập nhật quy trình nghiệp vụ sẽ được đáp ứng vào hệ thống dễ dàng.\\

Bằng cách sử dụng các công cụ và công nghệ đã trình bày ở trên, nhóm sẽ thực hiện tự động hóa quy trình nghiệp vụ doanh nghiệp, và áp dụng vào một hệ thống cửa hàng thời trang trong đề tài: "Xây dựng ứng dụng quản lý hệ thống cửa hàng thời trang dựa trên việc tự động hóa quy trình nghiệp vụ".
%%%%%%%%%%%%%%%%%%%%%%%%%

\section{Mục tiêu đề tài}
\hspace{0.5cm}Mục tiêu của đề tài là nghiên cứu và áp dụng tự động hóa quy trình nghiệp vụ vào trong một hệ thống cửa hàng thời trang. Việc vận hành hệ thống cần đảm bảo các yêu cầu chức năng và phi chức năng, đồng thời bao gồm các tính năng đã đề ra. Về phần quy trình nghiệp vụ, nghiệp vụ của hệ thống được quản lý và tự động hóa bằng phương pháp tự động hóa quy trình nghiệp vụ. Việc thiết kế và cập nhật nghiệp vụ sẽ được diễn ra nhanh chóng, dễ dàng dưới sự trợ giúp của BPMN và BPEL.

%%%%%%%%%%%%%%%%%%%%%%%%%
\section{Phạm vi đề tài}
\hspace{0.5cm}Trong đề tài này, nhóm sẽ áp dụng BPMN để thiết kế các quy trình nghiệp vụ cho một hệ thống cửa hàng thời trang, và sử dụng BPEL để thực thi quy trình nghiệp vụ. Nhóm sẽ phân tích, thiết kế và hiện thực một hệ thống bán hàng thời trang, bao gồm các tính năng mua / bán hàng trực tuyến và trực tiếp, và các tính năng quản lý hoạt động của hệ thống bán hàng.

%%%%%%%%%%%%%%%%%%%%%%%%%
\section{Ý nghĩa đề tài}
\hspace{0.5cm} Trong thực tế, quy trình nghiệp vụ của doanh nghiệp thường lớn và phức tạp, đòi hỏi sự tham gia của nhiều thành phần có các mối liên hệ phức tạp với nhau. Một sự thay đổi nghiệp vụ đòi hỏi nhiều thời gian và công sức để áp dụng vào hệ thống thực tế. Đối với phần lập trình và phát triển hệ thống, lập trình viên cần thay đổi mã nguồn ở nhiều nơi để đáp ứng đúng với thay đổi của nghiệp vụ, và cần đảm bảo sau khi thay đổi hệ thống vừa chạy đúng với nghiệp vụ mới, vừa đáp ứng các yêu cầu chức năng và phi chức năng đã đặt ra từ đầu. Với cách tiếp cận này, việc thay đổi và cập nhật quy trình nghiệp vụ sẽ tốn nhiều thời gian khi phải thay đổi mã nguồn ở nhiều phần trong hệ thống, sẽ tốn công sức để đảm bảo hệ thống đáp ứng đúng như yêu cầu sau cập nhật, và dễ có rủi ro nếu việc tìm và thay đổi mã nguồn vẫn chưa triệt để.\\

Với sự giúp đỡ của việc tự động hóa quy trình nghiệp vụ, việc quản lý quy trình nghiệp vụ sẽ trở nên dễ dàng hơn rất nhiều. Quy trình nghiệp vụ có thể được biểu diễn dưới dạng mô hình BPMN - một mô hình có thể hiểu được bởi người dùng nghiệp vụ và lập trình viên. Việc thay đổi và cập nhật quy trình nghiệp vụ sẽ được tối ưu và tự động với Ngôn ngữ thực thi quy trình nghiệp vụ - BPEL. Các thành phần của hệ thống tương ứng với các bước trong quy trình nghiệp vụ sẽ được kết nối đến các thành phần của BPEL, do đó khi có sự thay đổi của quy trình nghiệp vụ, chỉ cần cập nhật lại BPEL thì BPEL sẽ gọi thứ tự thực thi nghiệp vụ mới theo đúng như nó được cập nhật. Do đó, lập trình viên không cần tìm và sửa mã ở khắp nơi trong mã nguồn, mà chỉ cần cập nhật BPEL và đảm bảo kết nối đúng các thành phần BPEL đến các thành phần trong hệ thống. Do BPEL là một ngôn ngữ thực thi, nó có thể được deploy và chạy, nên khi có sự thay đổi nghiệp vụ ở BPMN và BPEL, việc cập nhật nghiệp vụ sẽ nhanh và tự động mà không yêu cầu nhiều công sức của con người.

%%%%%%%%%%%%%%%%%%%%%%%%%
\section{Đóng góp của đề tài}
\hspace{0.5cm} Đề tài mong muốn chứng minh sự hiệu quả khi áp dụng BPMN và BPEL vào việc thiết kế và quản lý quy trình trong mỗi doanh nghiệp. Sau khi hoàn thành, đề tài muốn cho thấy cách tiếp cận bằng tự động hóa quy trình nghiệp vụ sẽ giảm đáng kể thời gian và nguồn lực trong việc quản lý quy trình nghiệp vụ, và sẽ mang lại hiệu quả lớn cho doanh nghiệp trong việc vận hành hệ thống của mình.

%%%%%%%%%%%%%%%%%%%%%%%%%   
\section{Phân chia công việc}
 {


  \begin{figure}[h]
	  \begin{center}
		  \includegraphics[width=14cm]{img/plan.png}
	  \end{center}
	  \caption{Biểu đồ Gantt quá trình thực hiện đồ án tốt nghiệp}
  \end{figure}


  \newpage

  \setlength\extrarowheight{6pt}
  \begin{longtable}{| p{2cm} | p{2cm} | p{10cm} |}

	  \hline
	  \textbf{Bắt đầu} & \textbf{Hạn} & \textbf{Công việc}                                             \\
	  \hline
	  02/01/2023       & 08/01/2022   &
	  - Phong: Setup môi trường cho CustomerUI và hiện thực layout.
	  \newline
	  - Hiển: Setup môi trường cho AdminUI và hiện thực layout.
	  \newline
	  - Thọ: Setup môi trường cho backend, khởi tạo các service và đảm bảo giao tiếp giữa các service. \\
	  \hline
	  09/01/2023       & 16/09/2022   &
	  - Phong: Hiện thực Login page - Register page và Reset Password page.
	  \newline
	  - Hiển: Hiện thực Branch page - Branch detail page - Order page - Order detail page.
	  \newline
	  - Thọ: Hiện thực Customer service                                                                \\
	  \hline
	  30/01/2023       & 05/02/2022   &
	  - Phong: Hiện thực Home page - List Product page.
	  \newline
	  - Hiển: Hiện thực Event page - Event detail page - Goods page - Goods detail page, Goods Edit page.
	  \newline
	  - Thọ: Hiện thực Cart service.                                                                   \\
	  \hline
	  06/02/2023       & 12/02/2022   &
	  - Phong: Hiện thực Product Detail page - Cart page.
	  \newline
	  - Hiển: Hiện thực Staff page - Staff Detail page - Staff Request page.
	  \newline
	  - Thọ: Hiện thực Order service                                                                   \\
	  \hline
	  13/02/2023       & 19/02/2022   &
	  - Phong: Hiện thực Payment page - Manager Order page.
	  \newline
	  - Hiển: Hiện thực Account page, Account Detail page.
	  \newline
	  - Thọ: Hiện thực Account service                                                                 \\
	  \hline
	  20/02/2023       & 26/02/2022   &
	  - Phong: Hiện thực Order Detail page và Customer Info page.
	  \newline
	  - Hiển: Hiện thực Warehouse page - Statistic page.
	  \newline
	  - Thọ: Hiện thực Staff service.                                                                  \\
	  \hline
	  27/02/2023       & 05/03/2022   &
	  - Phong: Refactor code cho CustomerUI, hiện thực Event service.
	  \newline
	  - Hiển: Refactor code cho AdminUI, hiện thực Statistic service.
	  \newline
	  - Thọ: Hiện thực Branch service.                                                                 \\
	  \hline
	  06/03/2023       & 12/03/2022   &
	  - Hiển           \& Phong: Nghiên cứu, chuyển đổi BPMN sang BPEL
	  \newline
	  - Thọ: Hiện thực Warehouse service.                                                              \\
	  \hline
	  13/03/2023       & 19/03/2022   &
	  - Hiển           \& Phong: Nghiên cứu, chuyển đổi BPMN sang BPEL
	  \newline
	  - Thọ: Hiện thực Goods service.                                                                  \\
	  \hline
	  20/03/2023       & 26/03/2022   &
	  - Hiển           \& Phong: Nghiên cứu, chuyển đổi BPMN sang BPEL
	  \newline
	  - Thọ: Hiện thực Event service.                                                                  \\
	  \hline
	  27/03/2023       & 02/04/2022   &
	  - Hiển           \& Phong: Nghiên cứu, chuyển đổi BPMN sang BPEL
	  \newline
	  - Thọ: Hiện thực Statistic service.
	  \newline
	  - Tất cả: Viết báo cáo                                                                           \\
	  \hline
	  03/04/2023       & 16/04/2022   &
	  - Hiển           \& Phong: Nghiên cứu, chuyển đổi BPMN sang BPEL.
	  \newline
	  - Thọ: Hiện thực BFF và các adapter
	  \newline
	  - Tất cả: Viết báo cáo                                                                           \\
	  \hline
	  17/04/2023       & 23/04/2022   &
	  - Hiển           \& Phong: Nghiên cứu, chuyển đổi BPMN sang BPEL.
	  \newline
	  - Thọ: Hiện thực kết nối Kafka
	  \newline
	  - Tất cả: Viết báo cáo                                                                           \\
	  \hline
	  24/04/2023       & 30/04/2022   &
	  - Tất cả thành viên: Tổng hợp hệ thống.
	  \newline
	  - Tất cả: Viết báo cáo                                                                           \\
	  \hline
	  01/05/2023       & 07/05/2022   &
	  - Hiển           \& Phong: Viết báo cáo và kiểm thử hệ thống
	  \newline
	  - Thọ: Triển khai hệ thống                                                                       \\
	  \hline
	  08/05/2023       & 14/05/2022   &
	  - Tất cả thành viên: Viết báo cáo, kiểm thử hệ thống                                             \\
	  \hline
	  15/05/2023       & 29/05/2022   &
	  - Tất cả thành viên: Chuẩn bị cho phản biện và bảo vệ đề tài                                     \\
	  \hline
  \end{longtable}

 }

%%%%%%%%%%%%%%%%%%%%%%%%%%%%%%%%%%%%

% Cơ sở lý thuyết
\chapter{Cơ sở lý thuyết}
% % !TeX root = ..\main.tex
\section{Cơ sở lý thuyết}
\subsection{Lược đồ BPMN}
\subsubsection{Tổng quan về BPMN}

\hspace*{0.5cm}BPMN(Business Process Modeling Notation) là một lược đồ tập hợp các kí hiệu để mô hình hóa trực quan các quy trình nghiệp vụ xử lý . BPMN giúp các doanh nghiệp trực quan hóa các hoạt động và các luồng thông tin để thực hiện quy trình nghiệp vụ.
\begin{figure}[!htp]
	\centering
	\includegraphics[width=10cm]{img/theory/BPMN/BPMN_sample.png}
	\newline
	\caption{Ví dụ về lược đồ BPMN}
\end{figure}



\subsubsection{Lợi ích của BPMN}
\begin{itemize}
	\item Giúp các bên liên quan có thể hiểu rõ quy trình nghiệp vụ của hệ thống
	\item Thu hẹp khoảng cách giửa bộ phận thiết kế và bộ phận nghiệp vụ
	\item Dễ dàng mô tả các nghiệp vụ phức tạp
\end{itemize}

\subsubsection{Các thành phần của BPMN}
\subsubsubsection*{Swimlane}
Swimlane bao gồm Pool và Lane:
\begin{itemize}
	\item Pool: Đại diện cho một tổ chức, phòng ban, một vai trò hoặc một hệ thống nào đó.
	\item Lane: Đại diện các cá nhân riêng lẻ, người sẽ làm các hoạt động cụ thể.
\end{itemize}

\begin{figure}[!htp]
	\centering
	\includegraphics[width=10cm]{img/theory/BPMN/BPMN_swimlane.png}
	\newline
	\caption{Lược đồ BPMN bao gồm 1 pool(Hệ thống cứa hàng) và 2 lane (Nhân viên, Quản lý chi nhánh)}
\end{figure}

\subsubsubsection*{Activities}
Mô tả công việc trong quy trình, được kí hiệu bằng một kình chử nhật bo tròn 4 góc. Bao gồm 4 loại:
\begin{itemize}
	\item Task: Các việc nhỏ, đơn lẻ
	      \begin{figure}[!htp]
		      \begin{center}
			      \includegraphics[width=2cm]{img/theory/BPMN/Task.png}
		      \end{center}
		      \caption{Task}
	      \end{figure}
	\item Event Sub-Process: Một quy trình con nằm trong một quy trình lớn, chức nhiều task
	      \begin{figure}[!htp]
		      \begin{center}
			      \includegraphics[width=2cm]{img/theory/BPMN/Event Sub-process.png}
		      \end{center}
		      \caption{Event Sub-Process}
	      \end{figure}
	\item Transaction: Các giao dịch, bao gồm các task nhỏ khác logic với nhau
	      \begin{figure}[!htp]
		      \begin{center}
			      \includegraphics[width=2cm]{img/theory/BPMN/Transaction.png}
		      \end{center}
		      \caption{Transaction}
	      \end{figure}
	\item Call Activity: Gọi một quy trình khác đã được định nghĩa trong hệ thống, thay vì phải định nghĩa lại nhiều lần
	      \begin{figure}[!htp]

		      \begin{center}
			      \includegraphics[width=2cm]{img/theory/BPMN/Call Activity.png}
		      \end{center}
		      \caption{Call Activity}
	      \end{figure}
\end{itemize}
\begin{flushleft}
	Các hình ảnh được lấy từ \url{https://thinhnotes.com/chuyen-nghe-ba/giai-ngo-cac-ky-hieu-bpmn/}
\end{flushleft}


\subsubsubsection*{Events}
Sự kiện là một điều gì đó xảy ra và có thể tác động đến quy trình nghiệp vụ. Một sự kiện có thể là bên ngoài hoặc bên trong nội bộ. Miễn là nó có tác động và ảnh hưởng trực tiếp đến quy trình nghiệp vụ. Các sự kiện thường được kí hiệu bằng hình tròn, bên trong có các kí hiệu về loại hình kích hoạt. Bao gồm các loại:
\begin{itemize}
	\item Start: Các sự kiện mở đầu quy trình nghiệp vụ
	\item Intermediate: Các sự kiện xảy ra giữa các quy trình nghiệp vụ
	\item End: Các sự kiện kết thúc quy trình nghiệp vụ
\end{itemize}
\begin{figure}[!htp]
	\begin{center}
		\includegraphics[width=6cm]{img/theory/BPMN/Event.png}
	\end{center}
	\caption{Các loại sự kiện trong BPMN}
\end{figure}
\begin{flushleft}
	Hình ảnh được lấy từ \textit{\url{https://www.edrawsoft.com/what-is-bpmn.html?gclid=Cj0KCQiAvqGcBhCJARIsAFQ5ke5U9xsSqU24T_1MgVtsStQ3wV0NbkzXcIY-rDCU8JAet5rGQMc6PyIaAgdxEALw_wcB}}
\end{flushleft}



\subsubsubsection*{Gateways}
Các cổng có nhiệm vụ chính là kiểm soát dòng chảy của quy trình nghiệp vụ. Cổng thường kí hiệu hình thoi, bên trong là kí hiệu các loại.
\begin{itemize}
	\item Start: Các sự kiện mở đầu quy trình nghiệp vụ
	\item Intermediate: Các sự kiện xảy ra giữa các quy trình nghiệp vụ
	\item End: Các sự kiện kết thúc quy trình nghiệp vụ
\end{itemize}
\begin{figure}[!htp]
	\begin{center}
		\includegraphics[width=8cm]{img/theory/BPMN/Gateway.png}
	\end{center}
	\caption{Các loại cổng phổ biến trong BPMN}
\end{figure}
\begin{flushleft}
	Hình ảnh được lấy từ \url{https://www.bacs.vn/vi/blog/kien-thuc/tong-quan-ve-bpmn-danh-cho-business-analyst-11002.html}
\end{flushleft}



%%%%%%%%%%%%%%%%%%%%%%%%%%%%%%%%%
\subsection{Ngôn ngữ BPEL}
\subsubsection{Tổng quan về BPEL}
BPEL hay BPEL4WS(Business Process Execution Language for Web Services) là ngôn ngữ dùng để định nghĩa và thực thi các quy trình nghiệp vụ sử dụng các dịch vụ web. Ngôn ngữ BPEL được xây dựng dựa trên nền tảng XML hỗ trợ các dịch vụ web: SOAP, WSDL, UDDI. BPEL xây dựng một dịch vụ web dựa trên việc điều phối, phối hợp các dịch vụ web độc lập khác để hình thành nên một dịch vụ web mới.
\subsubsection{Các thành phần trong BPEL}
BPEL bao gồm các thành phần nguyên tố và các thành phần cấu trúc. Các thành phần nguyên tố bao gồm:
\begin{itemize}
	\item $<invoke>$: Gọi một dịch vụ web khác
	\item $<receive>$: Đợi phản hồi tự các tác vụ bất đồng bộ
	\item $<reply>$: Phản hồi cho các tác vụ đồng bộ
	\item $<assign>$: Gán giá trị cho biến
	\item $<throw>$: Chỉ định các lỗi, ngoại lệ
	\item $<wait>$: Đợi một thời gian
	\item $<terminate>$: Kết thúc toàn bộ quy trình
\end{itemize}

Bên cạnh thành nguyên tố, các thành phần cấu trúc chứa nhiều tác vụ bên trong:
\begin{itemize}
	\item $<sequence>$: Các tác vụ thực hiện tuần tự
	\item $<flow>$: Các tác vụ thực hiện song song
	\item $<switch>$: Thực hiện các tác vụ rẽ nhánh
	\item $<while>$: Thực hiện các vọng lặp
\end{itemize}

Ngoài ra, còn có $<partnerLink>$ dùng để xác định các đối tác, các dịch vụ web sử dụng, $<variable>$ để định nghĩa các biến cần thiết trong quy trình

\subsubsection{Ví dụ về BPEL}


%%%%%%%%%%%%%%%%%%%%%%%%%%%%%%%%%
\subsection{Chuyển đổi BPMN sang BPEL}

%%%%%%%%%%%%%%%%%%%%%%%%%%%%%%%%%%%%

\chapter{Thiết kế hệ thống}
% Requirement
% \input{tex/Requirement}

\newpage
% Usecase
% % !TeX root = ..\main.tex
\section{Lược đồ Use-case}
\subsection{Toàn hệ thống}
\begin{figure}[!htp]
    \centering    \includegraphics[width=16.5cm]{img/UseCase/UseCase-full.drawio.png}
    \newline
    \caption{Use case Hệ thống quản lý cửa hàng thời trang}
\end{figure}
\textbf{Lưu ý:} Các Usecase không yêu cầu đăng nhập:
\begin{itemize}
    \item Đăng ký
    \item Tạo mới mật khẩu
    \item Xem danh sách sản phẩm
    \item Lọc sản phẩm
    \item Điểm danh
\end{itemize}
\newpage


\subsection{Đặt hàng trực tuyến}

\begin{figure}[!htp]
    \centering
    \includegraphics[width=12cm]{img/UseCase/UseCase-Đặt hàng.drawio.png}
    \newline
    \caption{Use case đặt hàng trực tuyến}
\end{figure}
\subsubsection{Xem danh sách sản phẩm}
\begin{center}
    {
        \setlength\extrarowheight{6pt}
        \begin{longtable}{| p{.20\textwidth} | p{.80\textwidth} |}
            \hline
            \textbf{Use-case name}
             &
            Xem danh sách sản phẩm
            \\
            \hline
            \textbf{Actor}
             &
            Người dùng
            \\
            \hline
            \textbf{Description}
             &
            Cho phép người dùng có thể xem danh sách các sản phẩm của cửa hàng.
            \\
            \hline
            \textbf{Preconditions}
             &
            Người dùng đã truy cập vào trang chủ của hệ thống.
            \\
            \hline
            \textbf{Postconditions}
             &
            Người dùng xem được danh sách các sản phẩm thời trang của cửa hàng.
            \\
            \hline
            \textbf{Trigger}
             &
            Không.
            \\
            \hline
            \textbf{Normal flow}
             &
            1. Hệ thống hiển thị các sản phẩm theo mặc định.

            2. Người dùng chọn mục tất cả sản phẩm.

            3. Hệ thống hiển thị toàn bộ sản phẩm có bán trong cửa hàng.

            4. Người dùng lướt trang web để duyệt tất cả các sản phẩm của cửa hàng.
            \\
            \hline
            \textbf{Alternative flow}
             &
            2.1 Người dùng lướt trang chủ để xem các sản phẩm nổi bật mặc định của trang web.
            \\
            %%%%%%%%%%%%%%%%
            \hline
            \textbf{Exceptions}
             &
            Không
            %%%%%%%%%%%%%%%%%%%%
            \\
            \hline
            \caption{Đặc tả use-case \textbf{xem danh sách sản phẩm}}
        \end{longtable}
    }

\end{center}


\subsubsection{Lọc sản phẩm theo từ khóa}
\begin{center}
    {
        \setlength\extrarowheight{6pt}
        \begin{longtable}{| p{.20\textwidth} | p{.80\textwidth} |}
            \hline
            \textbf{Use-case name}
             &
            Lọc danh sách sản phẩm theo từ khóa
            \\
            \hline
            \textbf{Actor}
             &
            Người dùng
            \\
            \hline
            \textbf{Description}
             &
            Cho phép người dùng có thể xem danh sách các sản phẩm của cửa hàng có chứa tên mà người dùng mong muốn.
            \\
            \hline
            \textbf{Preconditions}
             &
            Người dùng đã truy cập vào trang chủ của nhà hàng.
            \\
            \hline
            \textbf{Postconditions}
             &
            Người dùng xem được danh sách các sản phẩm  của cửa hàng có tên chứa từ khóa mà người dùng muốn xem.
            \\
            \hline
            \textbf{Trigger}
             &
            Không.
            \\
            \hline
            \textbf{Normal flow}
             &
            1. Hệ thống hiển thị các sản phẩm theo mặc định.

            2. Người dùng nhập từ khóa của sản phẩm muốn tìm.

            3. Hệ thống hiển thị toàn bộ sản phẩm mà tên có từ khóa mà khách hàng tìm kiếm.

            4. Người dùng xem kết quả danh sách sản phẩm tìm kiếm của hệ thống.
            \\
            \hline
            \textbf{Alternative flow}
             &
            Không
            %%%%%%%%%%%%%%%%
            \\
            \hline
            \textbf{Exceptions}
             &
            Không
            %%%%%%%%%%%%%%%%%%%%
            \\
            \hline
            \caption{Đặc tả use-case \textbf{lọc sản phẩm theo từ khóa}}
        \end{longtable}
    }
\end{center}


\subsubsection{Lọc sản phẩm theo sự kiện }
\begin{center}
    {
        \setlength\extrarowheight{6pt}
        \begin{longtable}{| p{.20\textwidth} | p{.80\textwidth} |}
            \hline
            \textbf{Use-case name}
             &
            Lọc danh sách sản phẩm theo sự kiện
            \\
            \hline
            \textbf{Actor}
             &
            Người dùng
            \\
            \hline
            \textbf{Description}
             &
            Cho phép người dùng có thể xem các sản phẩm của cửa hàng đang trong danh sách thuộc sự kiện mà người dùng muốn xem.
            \\
            \hline
            \textbf{Preconditions}
             &
            Người dùng đã truy cập vào trang chủ của nhà hàng.
            \\
            \hline
            \textbf{Postconditions}
             &
            Người dùng xem được các sản phẩm của cửa hàng đang trong danh sách thuộc sự kiện mà người dùng muốn xem.
            \\
            \hline
            \textbf{Trigger}
             &
            Không
            \\
            \hline
            \textbf{Normal flow}
             &
            1. Hệ thống hiển thị các sự kiện trong trang mặc định.

            2. Người dùng chọn sự kiện muốn xem.

            3. Hệ thống hiển thị toàn bộ sản phẩm mà có trong danh sách thuộc sự kiện người dùng chọn.

            4. Người dùng xem kết quả danh sách sản phẩm sau khi lọc.
            \\
            \hline
            \textbf{Alternative flow}
             &
            Không
            %%%%%%%%%%%%%%%%
            \\
            \hline
            \textbf{Exceptions}
             &
            Không
            %%%%%%%%%%%%%%%%%%%%
            \\
            \hline
            \caption{Đặc tả use-case \textbf{lọc sản phẩm theo sự kiện}}
        \end{longtable}
    }
\end{center}


\subsubsection{Lọc sản phẩm theo đặc điểm }
\begin{center}
    {
        \setlength\extrarowheight{6pt}
        \begin{longtable}{| p{.20\textwidth} | p{.80\textwidth} |}
            \hline
            \textbf{Use-case name}
             &
            Lọc danh sách sản phẩm theo đặc điểm
            \\
            \hline
            \textbf{Actor}
             &
            Người dùng
            \\
            \hline
            \textbf{Description}
             &
            Cho phép người dùng có thể xem các sản phẩm của cửa hàng có các đặc điểm của các thể loại mà người dùng chọn.
            \\
            \hline
            \textbf{Preconditions}
             &
            Người dùng đã vào trang tất cả sản phẩm của nhà hàng.
            \\
            \hline
            \textbf{Postconditions}
             &
            Người dùng xem được các sản phẩm của cửa hàng có các đặc điểm của thể loại mà người dùng muốn xem.
            \\
            \hline
            \textbf{Trigger}
             &
            Người dùng chọn "Tất cả sản phẩm".
            \\
            \hline
            \begin{flushleft}
                \textbf{Normal flow}
            \end{flushleft}
             &
            1. Người dùng chọn đặc điểm của các trường muốn lọc (Các trường này theo bao gồm giá, thương hiệu, gender, size, màu sắc, nhóm sản phẩm).

            2. Hệ thống hiển thị toàn bộ sản phẩm mà có các đặc điểm mà người dùng chọn.

            3. Người dùng xem kết quả danh sách sản phẩm sau khi lọc.
            \\
            \hline
            \textbf{Alternative flow}
             &
            Không
            %%%%%%%%%%%%%%%%
            \\
            \hline
            \textbf{Exceptions}
             &
            Không
            %%%%%%%%%%%%%%%%%%%%
            \\
            \hline
            \caption{Đặc tả use-case \textbf{lọc sản phẩm theo đặc điểm}}
        \end{longtable}
    }

\end{center}


\subsubsection{Xem thông tin chi tiết sản phẩm }
{
    \setlength\extrarowheight{6pt}
    \begin{longtable}{| p{.20\textwidth} | p{.80\textwidth} |}
        \hline
        \textbf{Use-case name}
         &
        Xem thông tin chi tiết sản phẩm
        \\
        \hline
        \textbf{Actor}
         &
        Người dùng
        \\
        \hline
        \textbf{Description}
         &
        Cho phép người dùng có thể xem thông tin chi tiết của sản phẩm.
        \\
        \hline
        \textbf{Preconditions}
         &
        Người dùng đang xem danh sách các sản phẩm.
        \\
        \hline
        \textbf{Postconditions}
         &
        Người dùng xem được thông tin chi tiết của sản phẩm mà người dùng muốn xem.
        \\
        \hline
        \begin{flushleft}
            \textbf{Normal flow}
        \end{flushleft}
         &
        1. Người dùng chọn vào sản phẩm muốn xem trong danh sách các sản phẩm đang xem.

        2. Hệ thống hiển thị thông tin chi tiết của sản phẩm.

        3. Người dùng xem thông tin chi tiết của sản phẩm.
        \\
        \hline
        \textbf{Alternative flow}
         &
        1.1. Người dùng xem danh sách sản phẩm trong giỏ hàng.

        1.2. Người dùng xem danh sách sản phẩm trong lịch sử mua hàng.
        %%%%%%%%%%%%%%%%
        \\
        \hline
        \textbf{Exceptions}
         &
        Không
        %%%%%%%%%%%%%%%%%%%%
        \\
        \hline
        \caption{Đặc tả use-case \textbf{xem thông tin chi tiết sản phẩm}}
    \end{longtable}
}


\subsubsection{ Thêm sản phẩm vào giỏ hàng}
{
    \setlength\extrarowheight{6pt}
    \begin{longtable}{| p{.20\textwidth} | p{.80\textwidth} |}
        \hline
        \textbf{Use-case name}
         &
        Thêm sản phẩm vào giỏ hàng
        \\
        \hline
        \textbf{Actor}
         &
        Người dùng
        \\
        \hline
        \textbf{Description}
         &
        Cho phép người dùng thêm sản phẩm vào giỏ hàng của tài khoản của họ.
        \\
        \hline
        \textbf{Preconditions}
         &
        Người dùng đã đăng nhập vào hệ thống.
        \\
        \hline
        \textbf{Postconditions}
         &
        Người dùng thêm sản phẩm vào giỏ hàng của tài khoản của họ.
        \\
        \hline
        \textbf{Triggers}
         &
        Không.
        \\
        \hline
        \begin{flushleft}
            \textbf{Normal flow}
        \end{flushleft}
         &
        1. Người dùng chọn vào xem danh sách các sản phẩm.

        2. Hệ thống hiển thị danh sách các sản phẩm của người dùng chọn

        3. Người dùng chọn xem thông tin chi tiết sản phẩm.

        4. Hệ thống hiển thị thông tin chi tiết sản phẩm.

        5. Người dùng chọn các thông tin sản phẩm theo ý muốn (Thông tin bao gồm size, màu... tùy thuộc loại sản phẩm).

        6. Hệ thống thực hiện yêu cầu và hiển thị kết quả.

        7. Người dùng nhận kết quả thêm vào giỏ hàng thành công.
        \\
        \hline
        \begin{flushleft}
            \textbf{Alternative flow}
        \end{flushleft}
         &
        1.1 Người dùng thêm vào giỏ từ lịch sử mua hàng.
        \begin{quote}
            1.1.1. Người dùng chọn xem đơn hàng đã mua.

            1.1.2. Hệ thống hiển thị danh sách các đơn hàng đã mua của người dùng.

            1.1.3. Người dùng chọn thêm lại vào giỏ hàng.

            1.1.4. Hệ thống thực hiện yêu cầu và hiển thị kết quả.

            1.1.5. Người dùng nhận kết quả thêm vào giỏ hàng thành công.

        \end{quote}
        3.1 Người dùng chọn thêm nhanh vào giỏ hàng.
        \begin{quote}
            3.1.1 Hệ thống thêm sản phẩm vào giỏ hàng của khách hàng với các thông tin mặc định của hệ thống và thông báo kết quả.

            3.1.2. Đến bước 7.
        \end{quote}
        %%%%%%%%%%%%%%%%
        \\
        \hline
        \begin{flushleft}
            \textbf{Exceptions}
        \end{flushleft}
         &
        7.1.Thêm vào giỏ hàng thất bại.
        \begin{quote}
            7.1.1 Quay lại bước 5.
        \end{quote}
        %%%%%%%%%%%%%%%%%%%%
        \\
        \hline
        \caption{Đặc tả use-case \textbf{thêm sản phẩm vào giỏ hàng}}
    \end{longtable}
}


\subsubsection{Mua hàng trực tuyến }
{
    \setlength\extrarowheight{6pt}
    \begin{longtable}{| p{.20\textwidth} | p{.80\textwidth} |}
        \hline
        \textbf{Use-case name}
         &
        Mua hàng trực tuyến
        \\
        \hline
        \textbf{Actor}
         &
        Người dùng
        \\
        \hline
        \textbf{Description}
         &
        Cho phép người dùng đặt hàng những sản phẩm mà họ muốn.
        \\
        \hline
        \textbf{Preconditions}
         &
        Người dùng đã đăng nhập vào hệ thống.
        \\
        \hline
        \textbf{Postconditions}
         &
        Người dùng đặt hàng thành công các sản phẩm mà họ muốn.
        \\
        \hline
        \begin{flushleft}
            \textbf{Normal flow}
        \end{flushleft}
         &
        1. Người dùng chọn giỏ hàng.

        2. Hệ thống hiển thị danh sách thông tin các sản phẩm trong giỏ hàng.

        3. Người dùng chỉnh sửa các thông tin sản phẩm theo ý muốn (Thông tin bao gồm size, màu... tùy thuộc loại sản phẩm).

        4. Hệ thống tính toán và hiển thị tổng giá tiền.

        5. Người dùng chọn đặt hàng.

        6. Hệ thống hiển thị màn hình danh sách thông tin sản phẩm mà người dùng đã chọn cùng giá tiền.

        7. Người dùng chọn và nhập địa chỉ giao hàng.

        8. Hệ thống nhận thông tin giao hàng và tính toán sau đó hiển thị kết quả lên màn hình.

        % 9. Người dùng nhập mã giảm giá.

        % 10. Hệ thống tính toán giá trị mã giảm giá và hiển thị giá cuối cùng cho người dùng.

        9. Người dùng chọn thanh toán và thực hiện thanh toán tùy chọn.

        10. Hệ thống thông báo kết quả đơn hàng hoàn tất cho người dùng.
        \\
        \hline
        \begin{flushleft}
            \textbf{Alternative flow}
        \end{flushleft}
         &
        3.1. Chưa có sản phẩm trong giỏ hàng.
        \begin{quote}
            3.1.1. Người dùng thực hiện tìm kiếm sản phẩm mình muốn và thêm vào giỏ hàng.

            3.1.2. Quay lại bước 1.
        \end{quote}
        5.1 Người dùng không muốn tiếp tục đặt hàng.
        \begin{quote}
            5.1.1. Người dùng chọn logo của trang web trên header.
            5.1.2. Hệ thống chuyển hướng về lại trang mặc định.
        \end{quote}
        % 9.1. Người dùng không muốn nhập mã giảm giá.
        % \begin{quote}
        %     9.1.1. Đến bước 11.
        % \end{quote}
        %%%%%%%%%%%%%%%%
        \\
        \hline
        \begin{flushleft}
            \textbf{Exceptions}
        \end{flushleft}
         &
        % 8.2. Địa chỉ khách hàng không được đơn vị vận chuyển hỗ trợ.
        % \begin{quote}
        %     8.2.1. Hệ thống thông báo cho người dùng và yêu cầu chỉnh sửa địa chỉ giao hàng.
        %     8.2.2. Quay lại bước 7.
        % \end{quote}
        % 10.1. Khách hàng thực hiện thanh toán thất bại khi chọn thanh toán online.
        % \begin{quote}
        %     10.1.1. Hệ thống thông báo cho người dùng đặt hàng thất bại và thực hiện lại thanh toán.
        %     10.1.2. Quay lại bước 9.
        % \end{quote}
        %%%%%%%%%%%%%%%%%%%%
        \\
        \hline
        \caption{Đặc tả use-case \textbf{mua hàng trực tuyến}}
    \end{longtable}
}
\subsubsection{Thanh toán}
{
    \setlength\extrarowheight{6pt}
    \begin{longtable}{| p{.20\textwidth} | p{.80\textwidth} |}
        \hline
        \textbf{Use-case name}
         &
        Thanh toán
        \\
        \hline
        \textbf{Actor}
         &
        Người dùng
        \\
        \hline
        \textbf{Description}
         &
        Cho phép người dùng có thể thực hiện thanh toán để hoàn tất đơn hàng.
        \\
        \hline
        \textbf{Preconditions}
         &
        Người dùng tiến hành đặt hàng.
        \\
        \hline
        \textbf{Postconditions}
         &
        Người dùng thanh toán đơn hàng thành công.
        \\
        \hline
        \textbf{Trigger}
         &
        Người dùng chọn "Thanh toán".
        \\
        \hline
        \begin{flushleft}
            \textbf{Normal flow}
        \end{flushleft}
         &
        1. Hệ thống hiển thị lựa chọn thanh toán cho khách hàng gồm "Thanh toán online" và "Thanh toán khi nhận hàng".

        2. Người dùng chọn "Thanh toán online".

        3. Hệ thống hiển thị các dịch vụ hỗ trợ thanh toán trực tuyến.

        4. Người dùng chọn dịch vụ thanh toán trực tuyến mà người dùng muốn.

        5. Hệ thống tổng hợp dữ liệu gửi cho bên dịch vụ thứ ba mà người dùng chọn sau đó liên kết hiển thị cho người dùng thực hiện thanh toán.

        6. Người dùng thực hiện thanh toán thông qua dịch vụ của bên thứ ba.

        7. Hệ thống bên thứ ba trả về kết quả thanh toán thành công cho hệ thống và hệ thống thực thông báo kết quả thanh toán cho người dùng.

        8. Người dùng xem kết quả giao dịch thanh toán.
        \\
        \hline
        \begin{flushleft}
            \textbf{Alternative flow}
        \end{flushleft}
         &
        2.1. Người dùng chọn "Thanh toán khi nhận hàng".
        \begin{quote}
            2.1.1. Hệ thống hiển thị xác nhận, ghi nợ cho vận chuyển và thông báo kết quả thành công.

            2.1.2. Người dùng xem kết quả.
        \end{quote}
        %%%%%%%%%%%%%%%%
        \\
        \hline
        \begin{flushleft}
            \textbf{Exceptions}
        \end{flushleft}
         &
        7.1. Thanh toán thất bại.
        \begin{quote}
            7.1.1. Hệ thống thông báo thanh toán thất bại.

            7.1.2. Quay lại bước 8.
        \end{quote}
        2.1.1.1. Xác nhận đơn hàng thất bại.
        \begin{quote}
            2.1.1.1.1 Hệ thống thông báo kết quả tạo đơn thanh toán khi nhận hàng thất bại.

            2.1.1.1.2 Quay lại bước 2.1.2.
        \end{quote}
        %%%%%%%%%%%%%%%%%%%%
        \\
        \hline
        \caption{Đặc tả use-case \textbf{Thanh toán}}
    \end{longtable}
}


\subsection{Quản lý sự kiện}
\newpage
\begin{figure}[!htp]
    \centering
    \includegraphics[width=8cm]{img/UseCase/UseCase-Quản lý sự kiện.drawio.png}
    \newline
    \caption{Use case quản lý sự kiện}
\end{figure}
\subsubsection{Xem thông tin chi tiết sự kiện}
{
    \setlength\extrarowheight{6pt}
    \begin{longtable}{| p{.20\textwidth} | p{.80\textwidth} |}
        \hline
        \textbf{Use-case name}
         &
        Xem thông tin chi tiết sự kiện.
        \\
        \hline
        \textbf{Actor}
         &
        Quản trị viên
        \\
        \hline
        \textbf{Description}
         &
        Cho phép quản trị viên xem  thông tin chi tiết sự kiện trong hệ thống.
        \\
        \hline
        \textbf{Preconditions}
         &
        Quản trị viên đã đăng nhập thành công vào hệ thống.
        \\
        \hline
        \textbf{Postconditions}
         &
        Quản trị viên xem được được thông tin chi tiết của sự kiện mà họ muốn.
        \\
        \hline
        \textbf{Trigger}
         &
        Quản trị viên chọn mục "sự kiện".
        \\
        \hline
        \begin{flushleft}
            \textbf{Normal flow}
        \end{flushleft}
         &
        1. Toàn bộ sự kiện được hiển thị lên trình duyệt.

        2.  Quản trị viên chọn sự kiện muốn xem thông tin chi tiết và nhấp chuột vào sự kiện đó.

        3.  Hệ thống hiển thị thông tin chi tiết của sự kiện.

        4. Người dùng xem thông tin chi tiết của sự kiện.
        \\
        \hline
        \textbf{Alternative flow}
         &
        Không.
        %%%%%%%%%%%%%%%%
        \\
        \hline
        \textbf{Exceptions}
         &
        Không.
        %%%%%%%%%%%%%%%%%%%%
        \\
        \hline
        \caption{Đặc tả use-case \textbf{xem thông tin chi tiết sự kiện}}
    \end{longtable}
}
\subsubsection{Thêm sự kiện cho hệ thống bán hàng online.}
{
    \setlength\extrarowheight{6pt}
    \begin{longtable}{| p{.20\textwidth} | p{.80\textwidth} |}
        \hline
        \textbf{Use-case name}
         &
        Thêm sự kiện cho hệ thống bán hàng online
        \\
        \hline
        \textbf{Actor}
         &
        Quản trị viên
        \\
        \hline
        \textbf{Description}
         &
        Cho phép quản trị viên có thể thêm sự kiện mới cho hệ thống bán hàng online.
        \\
        \hline
        \textbf{Preconditions}
         &
        Quản trị viên đã đăng nhập thành công vào hệ thống.
        \\
        \hline
        \textbf{Postconditions}
         &
        Quản trị viên tạo mới được sự kiện mà họ muốn.
        \\
        \hline
        \textbf{Trigger}
         &
        Quản trị viên chọn mục "sự kiện".
        \\
        \hline
        \begin{flushleft}
            \textbf{Normal flow}
        \end{flushleft}
         &
        1. Toàn bộ sự kiện được hiển thị lên trình duyệt.

        2. Quản trị viên chọn nút thêm sự kiện.

        3. Hệ thống hiển thị màn hình form để tạo sự kiện mới.

        4. Quản trị viên nhập các thông tin cần thiết mà hệ thống yêu cầu và chọn icon "Thêm sự kiện".

        5. Hệ thống cập nhật thêm sự kiện mới vào database.
        \\
        \hline
        \begin{flushleft}
            \textbf{Alternative flow}
        \end{flushleft}
         &
        4.1. Người dùng không muốn tiếp tục nữa.
        \begin{quote}
            4.1.1. Người dùng chọn "Hủy".

            4.1.2. Hệ thống tắt màn hình popup.
        \end{quote}
        %%%%%%%%%%%%%%%%
        \\
        \hline
        \begin{flushleft}
            \textbf{Exceptions}
        \end{flushleft}
         &
        Không
        %%%%%%%%%%%%%%%%%%%%
        \\
        \hline
        \caption{Đặc tả use-case \textbf{thêm sự kiện cho hệ thống bán hàng online}}
    \end{longtable}
}

\subsubsection{Chỉnh sửa sự kiện}
{
    \setlength\extrarowheight{6pt}
    \begin{longtable}{| p{.20\textwidth} | p{.80\textwidth} |}
        \hline
        \textbf{Use-case name}
         &
        Chỉnh sửa thông tin sự kiện
        \\
        \hline
        \textbf{Actor}
         &
        Quản trị viên
        \\
        \hline
        \textbf{Description}
         &
        Cho phép quản trị viên có thể chỉnh sửa thông tin sự kiện mới cho hệ thống bán hàng online.
        \\
        \hline
        \textbf{Preconditions}
         &
        Quản trị viên đã đăng nhập thành công vào hệ thống.
        \\
        \hline
        \textbf{Postconditions}
         &
        Quản trị viên chỉnh sửa được sự kiện mà họ muốn.
        \\
        \hline
        \textbf{Trigger}
         &
        Quản trị viên chọn mục "sự kiện".
        \\
        \hline
        \begin{flushleft}
            \textbf{Normal flow}
        \end{flushleft}
         &
        1.Toàn bộ sự kiện được hiển thị lên trình duyệt.

        2. Quản trị viên chọn sự kiện muốn chỉnh sửa và nhấn vào nút "chỉnh sửa".

        3. Hệ thống hiển thị màn hình các thông tin của sự kiện.

        4. Quản trị viên chỉnh sửa các thông tin mà họ muốn sau đó nhấn nút "Lưu".

        5. Hệ thống thông báo kết quả cập nhật thêm sự kiện mới vào database.
        \\
        \hline
        \begin{flushleft}
            \textbf{Alternative flow}
        \end{flushleft}
         &
        4.1. Người dùng không muốn tiếp tục nữa.
        \begin{quote}
            4.1.1. Người dùng chọn "Hủy".

            4.1.2. Hệ thống tắt màn hình popup.
        \end{quote}
        %%%%%%%%%%%%%%%%
        \\
        \hline
        \begin{flushleft}
            \textbf{Exceptions}
        \end{flushleft}
         &
        5.1. Thông tin thay đổi không hợp lệ
        \begin{quote}
            5.1.1. Hệ thống thông báo chỉnh sửa sự kiện thất bại

            5.1.2. Trở về bước 4.
        \end{quote}
        %%%%%%%%%%%%%%%%%%%%
        \\
        \hline
        \caption{Đặc tả use-case \textbf{chỉnh sửa thông tin sự kiện}}
    \end{longtable}
}

\subsection{Quản lý chi nhánh}
\begin{figure}[!htp]
    \centering
    \includegraphics[width=9cm]{img/UseCase/UseCase-Quản lý chi nhánh.drawio.png}
    \newline
    \caption{Use case quản lý chi nhánh}
\end{figure}
\subsubsection{Thêm chi nhánh}
{
    \setlength\extrarowheight{6pt}
    \begin{longtable}{| p{.20\textwidth} | p{.80\textwidth} |}
        \hline
        \textbf{Use-case name}
         &
        Thêm chi nhánh
        \\
        %%%%%%%%%%%%%%%%
        \hline
        \textbf{Actor}
         &
        Quản lý chi nhánh
        \\
        %%%%%%%%%%%%%%%%
        \hline
        \textbf{Description}
         &
        Quản lý chi nhánh có thể thêm chi nhánh mới
        \\
        %%%%%%%%%%%%%%%%
        \hline
        \textbf{Preconditions}
         &
        Quản lý chi nhánh đã đăng nhập
        \\
        %%%%%%%%%%%%%%%%
        \hline
        \textbf{Postconditions}
         &
        Chi nhánh mới được cập nhật
        \\
        %%%%%%%%%%%%%%%%
        \hline
        \textbf{Trigger}
         &
        Quản lý chi nhánh chọn vào quản lý chi nhánh
        \\
        %%%%%%%%%%%%%%%%
        \hline
        \begin{flushleft}
            \textbf{Normal flow}
        \end{flushleft}
         &
        1. Hệ thống hiển thị danh sách chi nhánh

        2. Quản lý chi nhánh chọn thêm mới

        3. Hệ thống hiển thị biểu mẫu nhập thông tin

        4. Quản trị viên nhập thông tin chi nhánh mới

        5. Quản trị viên nhập chọn xác nhận

        6. Hệ thống cập nhật chi nhánh mới
        \\
        %%%%%%%%%%%%%%%%
        \hline
        \begin{flushleft}
            \textbf{Alternative flow}
        \end{flushleft}
         &
        5.1. Quản trị viên nhập chọn hủy
        \begin{quote}
            5.1.1. Kết thúc
        \end{quote}
        \\
        %%%%%%%%%%%%%%%%
        \hline
        \begin{flushleft}
            \textbf{Exceptions}
        \end{flushleft}
         &
        Không
        \\
        \hline
        %%%%%%%%%%%%%%%%
        \caption{Đặc tả use-case \textbf{thêm chi nhánh}}
    \end{longtable}
}
\subsubsection{Chỉnh sửa thông tin chi nhánh}
{
    \setlength\extrarowheight{6pt}
    \begin{longtable}{| p{.20\textwidth} | p{.80\textwidth} |}
        \hline
        \textbf{Use-case name}
         &
        Chỉnh sửa tin chi nhánh
        \\
        %%%%%%%%%%%%%%%%
        \hline
        \textbf{Actor}
         &
        Quản lý chi nhánh
        \\
        %%%%%%%%%%%%%%%%
        \hline
        \textbf{Description}
         &
        Quản lý chi nhánh có thể chỉnh sửa thông tin chi nhánh
        \\
        %%%%%%%%%%%%%%%%
        \hline
        \textbf{Preconditions}
         &
        Quản lý chi nhánh đã đăng nhập
        \\
        %%%%%%%%%%%%%%%%
        \hline
        \textbf{Postconditions}
         &
        Thông tin chi nhánh được cập nhật
        \\
        %%%%%%%%%%%%%%%%
        \hline
        \textbf{Trigger}
         &
        Quản lý chi nhánh chọn vào quản lý chi nhánh
        \\
        %%%%%%%%%%%%%%%%
        \hline
        \begin{flushleft}
            \textbf{Normal flow}
        \end{flushleft}
         &
        1. Hệ thống hiển thị danh sách chi nhánh

        2. Quản lý chi nhánh tìm kiếm chi nhánh và chọn chỉnh sửa

        3. Hệ thống hiển thị biểu mẫu nhập thông tin

        4. Quản lý chi nhánh chỉnh sửa thông tin chi nhánh

        5. Quản lý chi nhánh nhập chọn xác nhận

        6. Hệ thống cập nhật thông tin  chi nhánh
        \\
        %%%%%%%%%%%%%%%%
        \hline
        \begin{flushleft}
            \textbf{Alternative flow}
        \end{flushleft}
         &
        5.1. Quản lý chi nhánh chọn hủy
        \begin{quote}
            5.1.1. Kết thúc
        \end{quote}
        \\
        %%%%%%%%%%%%%%%%
        \hline
        \begin{flushleft}
            \textbf{Exceptions}
        \end{flushleft}
         &
        6.1. Tên chi nhánh bị trùng
        \begin{quote}
            6.1.1. Quay về bước 3
        \end{quote}
        \\
        \hline
        %%%%%%%%%%%%%%%%
        \caption{Đặc tả use-case \textbf{chỉnh sửa chi nhánh}}
    \end{longtable}
}

\subsubsection{Xem thông tin chi tiết chi nhánh}
{
    \setlength\extrarowheight{6pt}
    \begin{longtable}{| p{.20\textwidth} | p{.80\textwidth} |}
        \hline
        \textbf{Use-case name}
         &
        Xem thông tin chi tiết chi nhánh
        \\
        %%%%%%%%%%%%%%%%
        \hline
        \textbf{Actor}
         &
        Quản lý chi nhánh
        \\
        %%%%%%%%%%%%%%%%
        \hline
        \textbf{Description}
         &
        Quản lý chi nhánh có thể xem thông tin chi nhánh
        \\
        %%%%%%%%%%%%%%%%
        \hline
        \textbf{Preconditions}
         &
        Quản lý chi nhánh đã đăng nhập
        \\
        %%%%%%%%%%%%%%%%
        \hline
        \textbf{Postconditions}
         &
        Nhận được thông tin chi nhánh
        \\
        %%%%%%%%%%%%%%%%
        \hline
        \textbf{Trigger}
         &
        Quản lý chi nhánh truy cập quản lý chi nhánh
        \\
        %%%%%%%%%%%%%%%%
        \hline
        \begin{flushleft}
            \textbf{Normal flow}
        \end{flushleft}
         &
        1. Hệ thống hiển thị danh sách các chi nhánh

        2. Quản lý chi nhánh tìm kiếm và chọn chi nhánh muốn xem

        3. Hệ thống hiển thị thông tin chi nhánh đã chọn
        \\
        %%%%%%%%%%%%%%%%
        \hline
        \textbf{Alternative flow}
         &
        Không
        \\
        %%%%%%%%%%%%%%%%
        \hline
        \textbf{Exceptions}
         &
        Không
        \\
        \hline
        %%%%%%%%%%%%%%%%
        \caption{Đặc tả use-case \textbf{xem thông tin chi tiết chi nhánh}}
    \end{longtable}
}
\subsubsection{Xóa chi nhánh}
{
    \setlength\extrarowheight{6pt}
    \begin{longtable}{| p{.20\textwidth} | p{.80\textwidth} |}
        \hline
        \textbf{Use-case name}
         &
        Xóa chi nhánh
        \\
        %%%%%%%%%%%%%%%%
        \hline
        \textbf{Actor}
         &
        Quản lý chi nhánh
        \\
        %%%%%%%%%%%%%%%%
        \hline
        \textbf{Description}
         &
        Quản lý chi nhánh có thể xóa chi nhánh
        \\
        %%%%%%%%%%%%%%%%
        \hline
        \textbf{Preconditions}
         &
        Quản lý chi nhánh đã đăng nhập
        \\
        %%%%%%%%%%%%%%%%
        \hline
        \textbf{Postconditions}
         &
        Chi nhánh bị xóa khỏi hệ thống
        \\
        %%%%%%%%%%%%%%%%
        \hline
        \textbf{Trigger}
         &
        Quản lý chi nhánh chọn quản lý chi nhánh
        \\
        %%%%%%%%%%%%%%%%
        \hline
        \begin{flushleft}
            \textbf{Normal flow}
        \end{flushleft}
         &
        1. Hệ thống hiển thị danh sách chi nhánh

        2. Quản lý chi nhánh tìm và chọn chi nhánh muốn xóa

        3. Hệ thống hiển thị thông tin chi tiết chi nhánh

        4. Quản lý chi nhánh chọn xóa chi nhánh

        5. Hệ thống yêu cầu xác nhận

        6. Quản lý chi nhánh chọn xác nhận

        7. Hệ thống cập nhật thông tin
        \\
        %%%%%%%%%%%%%%%%
        \hline
        \begin{flushleft}
            \textbf{Alternative flow}
        \end{flushleft}
         &
        6.1. Quản trị viên chọn hủy
        \begin{quote}
            6.1.1. Quay lại bước 3
        \end{quote}
        \\
        %%%%%%%%%%%%%%%%
        \hline
        \textbf{Exceptions}
         &
        Không
        \\
        \hline
        %%%%%%%%%%%%%%%%
        \caption{Đặc tả use-case \textbf{xóa chi nhánh}}
    \end{longtable}
}

\subsection{Quản lý hoạt động kinh doanh}
\begin{figure}[!htp]
    \centering
    \includegraphics[width=8cm]{img/UseCase/UseCase-Quản lý hoạt động kinh doanh.drawio.png}
    \newline
    \caption{Use case quản lý hoạt động kinh doanh}
\end{figure}
\subsubsection{Xem biểu đồ kinh doanh}
{
    \setlength\extrarowheight{6pt}
    \begin{longtable}{| p{.20\textwidth} | p{.80\textwidth} |}
        \hline
        \textbf{Use-case name}
         &
        Xem biểu đồ kinh doanh
        \\
        %%%%%%%%%%%%%%%%
        \hline
        \textbf{Actor}
         &
        Quản lý chi nhánh, trưởng chi nhánh
        \\
        %%%%%%%%%%%%%%%%
        \hline
        \textbf{Description}
         &
        Quản lý chi nhánh, trưởng chi nhánh có thể xem biểu đồ kinh doanh
        \\
        %%%%%%%%%%%%%%%%
        \hline
        \textbf{Preconditions}
         &
        Quản lý chi nhánh, trưởng chi nhánh đã đăng nhập
        \\
        %%%%%%%%%%%%%%%%
        \hline
        \textbf{Postconditions}
         &
        Nhận được biểu đồ kinh doanh
        \\
        %%%%%%%%%%%%%%%%
        \hline
        \textbf{Trigger}
         &
        Quản lý chi nhánh, trưởng chi nhánh chọn quản lý hoạt động kinh doanh
        \\
        %%%%%%%%%%%%%%%%
        \hline
        \begin{flushleft}
            \textbf{Normal flow}
        \end{flushleft}
         &
        1. Quản lý chi nhánh, trưởng chi nhánh loại hàng muốn thống kê

        2. Quản lý chi nhánh, trưởng chi nhánh chọn khoảng thời gian

        3. Quản lý chi nhánh, trưởng chi nhánh chọn loại biểu đồ

        4. Hệ thống hiển thị biểu đồ
        \\
        %%%%%%%%%%%%%%%%
        \hline
        \textbf{Alternative flow}
         &
        Không
        \\
        %%%%%%%%%%%%%%%%
        \hline
        \textbf{Exceptions}
         &
        Không
        \\
        \hline
        %%%%%%%%%%%%%%%%
        \caption{Đặc tả use-case \textbf{xem biểu đồ kinh doanh}}
    \end{longtable}
}


\subsection{Quản lý nhân viên}
\begin{figure}[!htp]
    \begin{center}
        \includegraphics[width=10cm]{img/UseCase/UseCase-Quản lý nhân viên.drawio.png}
    \end{center}
    \caption{Use case quản lý nhân viên}
\end{figure}
\newpage


\subsubsection{Xem thông tin chi tiết nhân viên}
{
    \setlength\extrarowheight{6pt}
    \begin{longtable}{| p{.20\textwidth} | p{.80\textwidth} |}
        \hline
        \textbf{Use-case name}
         &
        Xem thông tin chi tiết nhân viên
        \\
        %%%%%%%%%%%%%%%%%%
        \hline
        \textbf{Actor}
         &
        Quản lý chi nhánh, trưởng chi nhánh
        \\
        %%%%%%%%%%%%%%%%%%
        \hline
        \textbf{Description}
         &
        Quản lý chi nhánh, trưởng chi nhánh có thể xem thông tin chi tiết nhân viên
        \\
        %%%%%%%%%%%%%%%%%%
        \hline
        \textbf{Preconditions}
         &
        Quản lý chi nhánh, trưởng chi nhánh đã đăng nhập
        \\
        %%%%%%%%%%%%%%%%%%
        \hline
        \textbf{Postconditions}
         &
        Nhận được thông tin chi tiết nhân viên
        \\
        %%%%%%%%%%%%%%%%%%
        \hline
        \textbf{Trigger}
         &
        Quản lý chi nhánh, trưởng chi nhánh chọn quản lý nhân viên
        \\
        %%%%%%%%%%%%%%%%%%
        \hline
        \begin{flushleft}
            \textbf{Normal flow}
        \end{flushleft}
         &
        1. Hệ thống hiển thị danh sách nhân viên

        2. Quản lý chi nhánh, trưởng chi nhánh tìm và chọn nhân viên muốn xem

        3. Hệ thống hiển thị thông tin chi tiết nhân viên
        \\
        %%%%%%%%%%%%%%%%%%
        \hline
        \textbf{Alternative flow}
         &
        Không
        \\
        %%%%%%%%%%%%%%%%%%
        \hline
        \textbf{Exceptions}
         &
        Không
        \\
        %%%%%%%%%%%%%%%%%%
        \hline
        \caption{Đặc tả use-case \textbf{xem thông tin chi tiết nhân viên}}
    \end{longtable}
}

\newpage

\subsubsection{Chỉnh sửa thông tin chi tiết nhân viên}
{
    \setlength\extrarowheight{6pt}
    \begin{longtable}{| p{.20\textwidth} | p{.80\textwidth} |}
        \hline
        \textbf{Use-case name}
         &
        Chỉnh sửa thông tin chi tiết nhân viên
        \\
        %%%%%%%%%%%%%%%%%%
        \hline
        \textbf{Actor}
         &
        Quản lý chi nhánh
        \\
        %%%%%%%%%%%%%%%%%%
        \hline
        \textbf{Description}
         &
        Quản lý chi nhánh có thể chỉnh sửa thông tin nhân viên.
        \\
        %%%%%%%%%%%%%%%%%%
        \hline
        \textbf{Preconditions}
         &
        Quản lý chi nhánh đã đăng nhập
        \\
        %%%%%%%%%%%%%%%%%%
        \hline
        \textbf{Postconditions}
         &
        Thông tin nhân viên được cập nhật
        \\
        %%%%%%%%%%%%%%%%%%
        \hline
        \textbf{Trigger}
         &
        Quản lý chi nhánh chọn quản lý nhân viên
        \\
        %%%%%%%%%%%%%%%%%%
        \hline
        \begin{flushleft}
            \textbf{Normal flow}
        \end{flushleft}
         &
        1. Hệ thống hiển thị danh sách  nhân viên

        2. Quản lý chi nhánh tìm và chọn nhân viên cần chỉnh sửa

        3. Hệ thống hiển thị biểu mẫu thông tin nhân viên

        4. Quản lý chỉnh sửa thông tin

        5. Quản lý chọn lưu

        6. Hệ thống cập nhật thông tin
        \\
        %%%%%%%%%%%%%%%%%%
        \hline
        \begin{flushleft}
            \textbf{Alternative flow}
        \end{flushleft}
         &
        5.1. Quản lý chọn hủy
        \begin{quote}
            5.1.1. Quay lại bước 3
        \end{quote}
        \\
        %%%%%%%%%%%%%%%%%%
        \hline
        \textbf{Exceptions}
         &
        Không
        \\
        %%%%%%%%%%%%%%%%%%
        \hline
        \caption{Đặc tả use-case \textbf{chỉnh sửa thông tin nhân viên}}
    \end{longtable}
}

\subsubsection{Thêm nhân viên}
{
    \setlength\extrarowheight{6pt}
    \begin{longtable}{| p{.20\textwidth} | p{.80\textwidth} |}
        \hline
        \textbf{Use-case name}
         &
        Thêm nhân viên
        \\
        %%%%%%%%%%%%%%%%%%
        \hline
        \textbf{Actor}
         &
        Quản lý chi nhánh
        \\
        %%%%%%%%%%%%%%%%%%
        \hline
        \textbf{Description}
         &
        Quản lý chi nhánh có thể thêm nhân viên
        \\
        %%%%%%%%%%%%%%%%%%
        \hline
        \textbf{Preconditions}
         &
        Quản lý chi nhánh đã đăng nhập
        \\
        %%%%%%%%%%%%%%%%%%
        \hline
        \textbf{Postconditions}
         &
        Thông tin nhân viên mới được cập nhật
        \\
        %%%%%%%%%%%%%%%%%%
        \hline
        \textbf{Trigger}
         &
        Quản lý chi nhánh chọn quản lý nhân viên
        \\
        %%%%%%%%%%%%%%%%%%
        \hline
        \begin{flushleft}
            \textbf{Normal flow}
        \end{flushleft}
         &
        1. Hệ thống hiển thị danh sách nhân viên

        2. Quản lý chi nhánh chọn thêm mới

        3. Hệ thống hiển thị biểu mẫu thông tin  nhân viên mới

        4. Quản lý chi nhánh nhập thông tin nhân viên

        5. Quản lý chi nhánh chọn lưu

        6. Hệ thống cập nhật thông tin
        \\
        %%%%%%%%%%%%%%%%%%
        \hline
        \begin{flushleft}
            \textbf{Alternative flow}
        \end{flushleft}
         &
        5.1. Quản lý chọn hủy
        \begin{quote}
            5.1.1. Quay lại bước 3
        \end{quote}
        \\
        %%%%%%%%%%%%%%%%%%
        \hline
        \textbf{Exceptions}
         &
        Không
        \\
        %%%%%%%%%%%%%%%%%%
        \hline
        \caption{Đặc tả use-case \textbf{Thêm nhân viên}}
    \end{longtable}
}

\subsubsection{Xóa nhân viên}
{
    \setlength\extrarowheight{6pt}
    \begin{longtable}{| p{.20\textwidth} | p{.80\textwidth} |}
        \hline
        \textbf{Use-case name}
         &
        Xóa nhân viên
        \\
        %%%%%%%%%%%%%%%%%%
        \hline
        \textbf{Actor}
         &
        Quản lý chi nhánh
        \\
        %%%%%%%%%%%%%%%%%%
        \hline
        \textbf{Description}
         &
        Quản lý chi nhánh có thể xóa nhân viên
        \\
        %%%%%%%%%%%%%%%%%%
        \hline
        \textbf{Preconditions}
         &
        Quản lý chi nhánh đã đăng nhập
        \\
        %%%%%%%%%%%%%%%%%%
        \hline
        \textbf{Postconditions}
         &
        Thông tin nhân viên được xóa
        \\
        %%%%%%%%%%%%%%%%%%
        \hline
        \textbf{Trigger}
         &
        Quản lý chi nhánh chọn quản lý nhân viên
        \\
        %%%%%%%%%%%%%%%%%%
        \hline
        \begin{flushleft}
            \textbf{Normal flow}
        \end{flushleft}
         &
        1. Hệ thống hiển thị danh sách nhân viên

        2. Quản lý chi nhánh tìm và chọn nhân viên muốn xóa

        3. Hệ thống hiển thị thông tin chi tiết nhân viên

        4. Quản lý chi nhánh chọn xóa

        5. Hệ thống yêu cầu xác nhận

        6. Quản lý chi nhánh chọn xác nhận

        7. Hệ thống cập nhật thông tin
        \\
        %%%%%%%%%%%%%%%%%%
        \hline
        \textbf{Alternative flow}
         &
        6.1. Quản lý chọn hủy
        \begin{quote}
            6.1.1. Quay lại bước 3
        \end{quote}
        \\
        %%%%%%%%%%%%%%%%%%
        \hline
        \textbf{Exceptions}
         &
        Không
        \\
        %%%%%%%%%%%%%%%%%%
        \hline
        \caption{Đặc tả use-case \textbf{Xóa nhân viên}}
    \end{longtable}
}


\subsubsection{Yêu cầu thêm nhân viên}
{
    \setlength\extrarowheight{6pt}
    \begin{longtable}{| p{.20\textwidth} | p{.80\textwidth} |}
        \hline
        \textbf{Use-case name}
         &
        Yêu cầu thêm nhân viên
        \\
        %%%%%%%%%%%%%%%%%%
        \hline
        \textbf{Actor}
         &
        Trưởng chi nhánh
        \\
        %%%%%%%%%%%%%%%%%%
        \hline
        \textbf{Description}
         &
        Trưởng chi nhánh có thể yêu cầu cấp trên thêm nhân viên.
        \\
        %%%%%%%%%%%%%%%%%%
        \hline
        \textbf{Preconditions}
         &
        Trưởng chi nhánh đã đăng nhập
        \\
        %%%%%%%%%%%%%%%%%%
        \hline
        \textbf{Postconditions}
         &
        Yêu cầu thêm được gửi đi
        \\
        %%%%%%%%%%%%%%%%%%
        \hline
        \textbf{Trigger}
         &
        Trưởng chi nhánh chọn vào quản lý nhân viên
        \\
        %%%%%%%%%%%%%%%%%%
        \hline
        \begin{flushleft}
            \textbf{Normal flow}
        \end{flushleft}
         &
        1. Hệ thống hiển thị danh sách nhân viên

        2. Trưởng chi nhánh chọn yêu cầu thêm mới

        3. Hệ thống hiển thị biểu mẫu nhập thông tin

        4. Trưởng chi nhánh nhập thông tin

        5. Trưởng chi nhánh chọn gửi

        6. Hệ thống gửi yêu cầu
        \\
        %%%%%%%%%%%%%%%%%%
        \hline
        \textbf{Alternative flow}
         &
        Không
        \\
        %%%%%%%%%%%%%%%%%%
        \hline
        Không
         &
        Không
        \\
        %%%%%%%%%%%%%%%%%%
        \hline
        \caption{Đặc tả use-case \textbf{Yêu cầu thêm nhân viên}}
    \end{longtable}
}

\subsubsection{Yêu cầu xóa nhân viên}
{
    \setlength\extrarowheight{6pt}
    \begin{longtable}{| p{.20\textwidth} | p{.80\textwidth} |}
        \hline
        \textbf{Use-case name}
         &
        Yêu cầu xóa nhân viên
        \\
        %%%%%%%%%%%%%%%%%%
        \hline
        \textbf{Actor}
         &
        Trưởng chi nhánh
        \\
        %%%%%%%%%%%%%%%%%%
        \hline
        \textbf{Description}
         &
        Trưởng chi nhánh có thể yêu cầu cấp trên xóa nhân viên.
        \\
        %%%%%%%%%%%%%%%%%%
        \hline
        \textbf{Preconditions}
         &
        Trưởng chi nhánh đã đăng nhập
        \\
        %%%%%%%%%%%%%%%%%%
        \hline
        \textbf{Postconditions}
         &
        Yêu cầu xóa được gửi đi
        \\
        %%%%%%%%%%%%%%%%%%
        \hline
        \textbf{Trigger}
         &
        Trưởng chi nhánh chọn quản lý nhân viên
        \\
        %%%%%%%%%%%%%%%%%%
        \hline
        \begin{flushleft}
            \textbf{Normal flow}
        \end{flushleft}
         &
        1. Hệ thống hiển thị danh sách nhân viên

        2. Trưởng chi nhánh tìm và chọn nhân viên muốn yêu cầu

        3. Hệ thống hiển thị thông tin nhân viên

        4. Trưởng chi nhánh chọn yêu cầu xóa

        5. Hệ thống gửi thông tin đi
        \\
        %%%%%%%%%%%%%%%%%%
        \hline
        \textbf{Alternative flow}
         &
        Không
        \\
        %%%%%%%%%%%%%%%%%%
        \hline
        \textbf{Exceptions}
         &
        Không
        \\
        %%%%%%%%%%%%%%%%%%
        \hline
        \caption{Đặc tả use-case \textbf{Yêu cầu xóa nhân viên}}
    \end{longtable}
}

\subsubsection{Duyệt yêu cầu cập nhật nhân viên}
{
    \setlength\extrarowheight{6pt}
    \begin{longtable}{| p{.20\textwidth} | p{.80\textwidth} |}
        \hline
        \textbf{Use-case name}
         &
        Duyệt yêu cầu thêm nhân viên
        \\
        %%%%%%%%%%%%%%%%%%
        \hline
        \textbf{Actor}
         &
        Quản lý chi nhánh
        \\
        %%%%%%%%%%%%%%%%%%
        \hline
        \textbf{Description}
         &
        Quản lý chi nhánh có thể thêm nhân viên từ yêu cầu
        \\
        %%%%%%%%%%%%%%%%%%
        \hline
        \textbf{Preconditions}
         &
        Quản lý chi nhánh đã đăng nhập
        \\
        %%%%%%%%%%%%%%%%%%
        \hline
        \textbf{Postconditions}
         &
        Thông tin nhân viên được thêm
        \\
        %%%%%%%%%%%%%%%%%%
        \hline
        \textbf{Trigger}
         &
        Quản lý chi nhánh chọn quản lý nhân viên
        \\
        %%%%%%%%%%%%%%%%%%
        \hline
        \begin{flushleft}
            \textbf{Normal flow}
        \end{flushleft}
         &
        1. Hệ thống hiển thị danh sách nhân viên

        2. Quản lý chi nhánh chọn được yêu cầu

        3. Hệ thống hiển thị danh sách yêu cầu nhân viên đang chờ

        4. Quản lý chi nhánh tìm và chọn duyệt nhân viên được yêu cầu

        5. Hệ thống yêu cầu xác nhận

        6. Quản lý chi nhánh chọn xác nhận

        7. Hệ thống cập nhật thông tin
        \\
        %%%%%%%%%%%%%%%%%%
        \hline
        \textbf{Alternative flow}
         &
        6.1. Quản lý chọn hủy
        \begin{quote}
            6.1.1. Quay lại bước 3
        \end{quote}
        \\
        %%%%%%%%%%%%%%%%%%
        \hline
        \textbf{Exceptions}
         &
        Không
        \\
        %%%%%%%%%%%%%%%%%%
        \hline
        \caption{Đặc tả use-case \textbf{Duyệt yêu cầu cập nhật nhân viên}}
    \end{longtable}
}


\newpage


\subsection{Quản lý tài khoản}
\begin{figure}[!htp]
    \centering
    \includegraphics[width=9cm]{img/UseCase/UseCase-Quản lý tài khoản.drawio.png}
    \newline
    \caption{Use case quản lý tài khoản}
\end{figure}

\subsubsection{Đăng ký}
{
    \setlength\extrarowheight{6pt}
    \begin{longtable}{| p{.20\textwidth} | p{.80\textwidth} |}
        \hline
        \textbf{Use-case name}
         &
        Đăng ký
        \\
        %%%%%%%%%%%%%%%%%%
        \hline
        \textbf{Actor}
         &
        Khách hàng
        \\
        %%%%%%%%%%%%%%%%%%
        \hline
        \textbf{Description}
         &
        Khách hàng có thể đăng ký tài khoản
        \\
        %%%%%%%%%%%%%%%%%%
        \hline
        \textbf{Preconditions}
         &
        Khách hàng truy cập hệ thống
        \\
        %%%%%%%%%%%%%%%%%%
        \hline
        \textbf{Postconditions}
         &
        Tài khoản được đăng ký
        \\
        %%%%%%%%%%%%%%%%%%
        \hline
        \textbf{Trigger}
         &
        Khách hàng chọn đăng ký
        \\
        %%%%%%%%%%%%%%%%%%
        \hline
        \begin{flushleft}
            \textbf{Normal flow}
        \end{flushleft}
         &
        1. Hệ thống hiển thị biểu mẫu thông tin khách hàng

        2. Khách hàng nhập thông tin của mình

        3. Khách hạng chọn đăng ký

        4. Hệ thống hiện yêu cầu nhập mã xác nhận

        5. Khách hàng nhập mã và chọn gửi

        6. Hệ thống cập nhật thông tin
        \\
        %%%%%%%%%%%%%%%%%%
        \hline
        \begin{flushleft}
            \textbf{Alternative flow}
        \end{flushleft}
         &
        3.1. Khách hàng chọn hủy
        \begin{quote}
            3.1.1 Kết thúc
        \end{quote}
        \\
        %%%%%%%%%%%%%%%%%%
        \hline
        \begin{flushleft}
            \textbf{Exceptions}
        \end{flushleft}
         &
        4.1. Số điện thoại đã được sử dụng
        \begin{quote}
            4.1.1 Quay lại bước 1
        \end{quote}
        6.1. Khách hàng nhập sai OTP
        \begin{quote}
            6.1.1 Quay lại bước 5
        \end{quote}
        \\
        %%%%%%%%%%%%%%%%%%
        \hline
        \caption{Đặc tả use-case \textbf{Đăng ký}}
    \end{longtable}
}

\subsubsection{Chỉnh sửa thông tin}
{
    \setlength\extrarowheight{6pt}
    \begin{longtable}{| p{.20\textwidth} | p{.80\textwidth} |}
        \hline
        \textbf{Use-case name}
         &
        Chỉnh sửa thông tin
        \\
        %%%%%%%%%%%%%%%%%%
        \hline
        \textbf{Actor}
         &
        Khách hàng
        \\
        %%%%%%%%%%%%%%%%%%
        \hline
        \textbf{Description}
         &
        Khách hàng có thể chỉnh sửa thông tin cá nhân
        \\
        %%%%%%%%%%%%%%%%%%
        \hline
        \textbf{Preconditions}
         &
        Khách hàng đã đăng nhập
        \\
        %%%%%%%%%%%%%%%%%%
        \hline
        \textbf{Postconditions}
         &
        Thông tin được cập nhật
        \\
        %%%%%%%%%%%%%%%%%%
        \hline
        \textbf{Trigger}
         &
        Khách hàng chọn vào tài khoản
        \\
        %%%%%%%%%%%%%%%%%%
        \hline
        \begin{flushleft}
            \textbf{Normal flow}
        \end{flushleft}
         &
        1. Hệ thống hiển thị biểu mẫu thông tin tài khoản

        2. Khách hàng chỉnh sửa thông tin tài khoản

        3. Khách hàng chọn lưu

        4. Hệ thống cập nhật thông tin
        \\
        %%%%%%%%%%%%%%%%%%
        \hline
        \begin{flushleft}
            \textbf{Alternative flow}
        \end{flushleft}
         &
        3.1.Khách hàng chọn hủy
        \begin{quote}
            3.1.1 Kết thúc
        \end{quote}
        \\
        %%%%%%%%%%%%%%%%%%
        \hline
        \textbf{Exceptions}
         &
        Không
        \\
        %%%%%%%%%%%%%%%%%%
        \hline
        \caption{Đặc tả use-case \textbf{Chỉnh sửa thông tin}}
    \end{longtable}
}

\subsubsection{Tạo mới mật khẩu}
{
    \setlength\extrarowheight{6pt}
    \begin{longtable}{| p{.20\textwidth} | p{.80\textwidth} |}
        \hline
        \textbf{Use-case name}
         &
        Tạo mới mật khẩu
        \\
        %%%%%%%%%%%%%%%%%%
        \hline
        \textbf{Actor}
         &
        Khách hàng
        \\
        %%%%%%%%%%%%%%%%%%
        \hline
        \textbf{Description}
         &
        Khách hàng có thể cập nhật mật khẩu mới khi quên mật khẩu
        \\
        %%%%%%%%%%%%%%%%%%
        \hline
        \textbf{Preconditions}
         &
        Khách hàng đã truy cập hệ thống
        \\
        %%%%%%%%%%%%%%%%%%
        \hline
        \textbf{Postconditions}
         &
        Thông tin được cập nhật
        \\
        %%%%%%%%%%%%%%%%%%
        \hline
        \textbf{Trigger}
         &
        Khách hàng chọn vào quên mật khẩu
        \\
        %%%%%%%%%%%%%%%%%%
        \hline
        \begin{flushleft}
            \textbf{Normal flow}
        \end{flushleft}
         &
        1. Hệ thống yêu cầu thông tin liên hệ

        2. Khách hàng nhập thông tin và gửi

        3. Hệ thống yêu cầu nhập mã xác nhận

        4. Khách hàng nhập mã và chọn gửi

        5. Hệ thống yêu cầu nhập mật khẩu mới

        6. Khách hàng nhập mật khẩu mới và chọn xác nhận

        7. Hệ thống cập nhật thông tin và hiển thị màn hình đăng nhập
        \\
        %%%%%%%%%%%%%%%%%%
        \hline
        \textbf{Alternative flow}
         &
        Không
        \\
        %%%%%%%%%%%%%%%%%%
        \hline
        \textbf{Exceptions}
         &
        6.1.  Khách hàng nhập sai OTP
        \begin{quote}
            6.1.1.Thông báo sai
            6.1.2 Quay lại bước 3
        \end{quote}
        \\
        %%%%%%%%%%%%%%%%%%
        \hline
        \caption{Đặc tả use-case \textbf{Tạo mới mật khẩu}}
    \end{longtable}
}

\subsubsection{Đổi mật khẩu}
{
    \setlength\extrarowheight{6pt}
    \begin{longtable}{| p{.20\textwidth} | p{.80\textwidth} |}
        \hline
        \textbf{Use-case name}
         &
        Đổi mật khẩu
        \\
        %%%%%%%%%%%%%%%%%%
        \hline
        \textbf{Actor}
         &
        Khách hàng, nhân viên
        \\
        %%%%%%%%%%%%%%%%%%
        \hline
        \textbf{Description}
         &
        Người dùng có thể đổi mật khẩu mới cho tài khoản
        \\
        %%%%%%%%%%%%%%%%%%
        \hline
        \textbf{Preconditions}
         &
        Người dùng đã đăng nhập hệ thống
        \\
        %%%%%%%%%%%%%%%%%%
        \hline
        \textbf{Postconditions}
         &
        Thông tin được cập nhật
        \\
        %%%%%%%%%%%%%%%%%%
        \hline
        \textbf{Trigger}
         &
        Khách hàng chọn vào đổi mật khẩu
        \\
        %%%%%%%%%%%%%%%%%%
        \hline
        \begin{flushleft}
            \textbf{Normal flow}
        \end{flushleft}
         &
        1. Hệ thống hiển thị form đổi mật khẩu

        2. Người dùng nhập thông tin và gửi

        3. Hệ thống kiểm tra thông tin, cập nhật và thông báo
        \\
        %%%%%%%%%%%%%%%%%%
        \hline
        \textbf{Alternative flow}
         &
        Không
        \\
        %%%%%%%%%%%%%%%%%%
        \hline
        \textbf{Exceptions}
         &
        Không
        \\
        %%%%%%%%%%%%%%%%%%
        \hline
        \caption{Đặc tả use-case \textbf{Đổi mật khẩu}}
    \end{longtable}
}



\subsection{Quản lý kho}
\begin{figure}[!htp]
    \centering
    \includegraphics[width=9cm]{img/UseCase/UseCase-Quản lý kho.drawio.png}
    \newline
    \caption{Use case quản lý kho}
\end{figure}


\subsubsection{Thêm kho}
{
    \setlength\extrarowheight{6pt}
    \begin{longtable}{| p{.20\textwidth} | p{.80\textwidth} |}
        \hline
        \textbf{Use-case name}
         &
        Thêm kho.
        \\
        \hline
        \textbf{Actor}
         &
        Quản lý kho.
        \\
        \hline
        \textbf{Description}
         &
        Quản lý thêm kho mới vào trong hệ thống lưu trữ.
        \\
        \hline
        \textbf{Preconditions}
         &
        Quản lý kho đã đăng nhập vào hệ thống.
        \\
        \hline
        \textbf{Postconditions}
         &
        Quản lý thêm được kho mà quản lý kho muốn.
        \\
        \hline
        \textbf{Trigger}
         &
        Quản lý kho chọn nút "Quản lý kho".
        \\
        \hline
        \begin{flushleft}
            \textbf{Normal flow}
        \end{flushleft}
         &
        1. Hệ thống hiển thị danh sách tất cả kho.
        2. Quản lý kho chọn nút "Thêm kho".

        3. Hệ thống hiển thị thông tin kho cần điền thông tin để thêm.

        4. Người dùng nhập thông tin của kho và nhấn "Tạo".

        5. Hệ thống tiến hành kiểm tra định dạng các thông tin mà người dùng chọn và thông báo kết quả tạo kho thành công.

        6. Người dùng xem kết quả tạo kho.
        \\
        \hline
        \textbf{Alternative flow}
         &
        Không.
        %%%%%%%%%%%%%%%%
        \\
        \hline
        \textbf{Exceptions}
         &
        5.1. Người dùng nhập thông tin không đúng định dạng.
        \begin{quote}
            5.1.1. Hệ thống thông báo thông tin kho nhập vào sai định dạng, yêu cầu người dùng nhập lại.

            5.1.2. Người dùng chỉnh sửa lại các thông tin theo định dạng và chọn "Tạo".

            5.1.3.Quay lại bước 5.
        \end{quote}
        5.2. Tạo kho thất bại.
        \begin{quote}
            5.2.1. Hệ thống thông báo kết quả tạo kho thất bại.
            5.2.2. Quay lại bước 2.
        \end{quote}
        %%%%%%%%%%%%%%%%%%%%
        \\
        \hline
        \caption{Đặc tả use-case \textbf{thêm kho}}
    \end{longtable}
}

\subsubsection{Chỉnh sửa thông tin kho}
{
    \setlength\extrarowheight{6pt}
    \begin{longtable}{| p{.20\textwidth} | p{.80\textwidth} |}
        \hline
        \textbf{Use-case name}
         &
        Chỉnh sửa thông tin kho.
        \\
        \hline
        \textbf{Actor}
         &
        Quản lý kho.
        \\
        \hline
        \textbf{Description}
         &
        Quản lý chỉnh sửa thông tin kho.
        \\
        \hline
        \textbf{Preconditions}
         &
        Quản lý kho đã đăng nhập vào hệ thống.
        \\
        \hline
        \textbf{Postconditions}
         &
        Quản lý chỉnh sửa được thông tin của kho mà quản lý kho muốn.
        \\
        \hline
        \textbf{Trigger}
         &
        Quản lý chọn nút "Quản lý kho".
        \\
        \hline
        \begin{flushleft}
            \textbf{Normal flow}
        \end{flushleft}
         &
        1. Hệ thống hiển thị danh sách tất cả kho.

        2. Quản lý kho chọn kho muốn chỉnh sửa thông tin và nhấn nút "Chỉnh sửa".

        3. Hệ thống hiển thị thông tin kho có sẵn.

        4. Người dùng chỉnh sửa thông tin của kho và nhấn "Lưu".

        5. Hệ thống tiến hành kiểm tra định dạng các thông tin mà người dùng chọn và thông báo kết quả chỉnh sửa kho thành công.

        6. Người dùng xem kết quả chỉnh sửa kho.
        \\
        \hline
        \begin{flushleft}
            \textbf{Alternative flow}
        \end{flushleft}
         &
        1.1. Quản lý kho muốn chỉnh sửa sau khi xem thông tin chi tiết của kho.
        \begin{quote}
            1.1.1. Quản lý kho chọn nút "Chỉnh sửa" trong trang xem thông tin chi tiết kho.
            1.1.2. Đến bước 3.
        \end{quote}
        %%%%%%%%%%%%%%%%
        \\
        \hline
        \begin{flushleft}
            \textbf{Exceptions}
        \end{flushleft}
         &
        5.1. Người dùng nhập thông tin không đúng định dạng.
        \begin{quote}
            5.1.1. Hệ thống thông báo thông tin kho nhập vào sai định dạng, yêu cầu người dùng nhập lại.
            5.1.2. Người dùng chỉnh sửa lại các thông tin theo định dạng và chọn "Lưu".
            5.1.3. Quay lại bước 5.
        \end{quote}
        5.2. Hệ thống chỉnh sửa kho thất bại.
        \begin{quote}
            5.2.1. Hệ thống thông báo kết quả chỉnh sửa kho thất bại.
            5.2.2. Quay lại bước 2.
        \end{quote}
        %%%%%%%%%%%%%%%%%%%%
        \\
        \hline
        \caption{Đặc tả use-case \textbf{chỉnh sửa thông tin kho}}
    \end{longtable}
}


\subsection{Quản lý hàng hóa}
\begin{figure}[!htp]
    \centering
    \includegraphics[width=10cm]{img/UseCase/UseCase-Quản lý hàng.drawio.png}
    \newline
    \caption{Use case quản lý hàng hóa}
\end{figure}


\subsubsection{Thêm hàng}
{
    \setlength\extrarowheight{6pt}
    \begin{longtable}{| p{.20\textwidth} | p{.80\textwidth} |}
        \hline
        \textbf{Use-case name}
         &
        Thêm hàng
        \\
        %%%%%%%%%%%%%%%%%%
        \hline
        \textbf{Actor}
         &
        Quản lý hàng hóa
        \\
        %%%%%%%%%%%%%%%%%%
        \hline
        \textbf{Description}
         &
        Quản lý hàng hóa có thể thêm vào hệ thống
        \\
        %%%%%%%%%%%%%%%%%%
        \hline
        \textbf{Preconditions}
         &
        Quản lý hàng hóa đã đăng nhập
        \\
        %%%%%%%%%%%%%%%%%%
        \hline
        \textbf{Postconditions}
         &
        Hàng được thêm vào hệ thống
        \\
        %%%%%%%%%%%%%%%%%%
        \hline
        \textbf{Trigger}
         &
        Quản lý hàng hóa chọn "Quản lý hàng hóa"
        \\
        %%%%%%%%%%%%%%%%%%
        \hline
        \begin{flushleft}
            \textbf{Normal flow}
        \end{flushleft}
         &
        1. Hệ thống hiển thị danh sách hàng hóa

        2. Quản lý chọn thêm mới

        3. Hệ thống hiển thị trang thêm mới hàng

        4. Quản lý hàng hóa nhập thông tin hàng

        5. Quản lý hàng hóa chọn "Lưu"

        6. Hệ thống cập nhật thông tin
        \\
        %%%%%%%%%%%%%%%%%%
        \hline
        \begin{flushleft}
            \textbf{Alternative flow}
        \end{flushleft}
         &
        5.1. Quản lý hàng hóa chọn "Hủy"
        \begin{quote}

            5.1.1 Quay lại bước 1
        \end{quote}
        \\
        %%%%%%%%%%%%%%%%%%
        \hline
        \textbf{Exceptions}
         &
        Không
        \\
        %%%%%%%%%%%%%%%%%%
        \hline
        \caption{Đặc tả use-case \textbf{Thêm hàng}}
    \end{longtable}
}

\subsubsection{Chỉnh sửa thông tin hàng}
{
    \setlength\extrarowheight{6pt}
    \begin{longtable}{| p{.20\textwidth} | p{.80\textwidth} |}
        \hline
        \textbf{Use-case name}
         &
        Chỉnh sửa thông tin hàng
        \\
        %%%%%%%%%%%%%%%%%%
        \hline
        \textbf{Actor}
         &
        Quản lý hàng hóa
        \\
        %%%%%%%%%%%%%%%%%%
        \hline
        \textbf{Description}
         &
        Quản lý hàng hóa có thể thông tin đầy đủ của hàng
        \\
        %%%%%%%%%%%%%%%%%%
        \hline
        \textbf{Preconditions}
         &
        Quản lý hàng hóa đã đăng nhập
        \\
        %%%%%%%%%%%%%%%%%%
        \hline
        \textbf{Postconditions}
         &
        Thông tin hàng được cập nhật
        \\
        %%%%%%%%%%%%%%%%%%
        \hline
        \textbf{Trigger}
         &
        Quản lý hàng hóa chọn "Quản lý hàng hóa"
        \\
        %%%%%%%%%%%%%%%%%%
        \hline
        \begin{flushleft}
            \textbf{Normal flow}
        \end{flushleft}
         &
        1. Hệ thống hiển thị danh sách hàng hóa

        2. Quản lý hàng hóa tìm kiếm và chọn hàng muốn chỉnh sửa

        3. Thông tin chi tiết hàng được hiển thị

        4. Quản lý hàng hóa chọn "Chỉnh sửa"

        5. Hệ thống hiển thị biểu mẫu thông tin hàng

        6. Quản lý hàng hóa nhập thông tin

        7. Quản lý hàng hóa chọn "Lưu"

        8. Hệ thống cập nhật thông tin
        \\
        %%%%%%%%%%%%%%%%%%
        \hline
        \begin{flushleft}
            \textbf{Alternative flow}
        \end{flushleft}
         &
        7.1 Quản lý hàng hóa chọn "Hủy"
        \begin{quote}
            7.1.1 Quay lại bước 3
        \end{quote}
        \\
        %%%%%%%%%%%%%%%%%%
        \hline
        \textbf{Exceptions}
         &
        Không
        \\
        %%%%%%%%%%%%%%%%%%
        \hline
        \caption{Đặc tả use-case \textbf{Chỉnh sửa thông tin chi tiết hàng}}
    \end{longtable}
}

\subsubsection{Thêm hàng vào kho}
{
    \setlength\extrarowheight{6pt}
    \begin{longtable}{| p{.20\textwidth} | p{.80\textwidth} |}
        \hline
        \textbf{Use-case name}
         &
        Thêm hàng vào kho
        \\
        %%%%%%%%%%%%%%%%%%
        \hline
        \textbf{Actor}
         &
        Quản lý hàng hóa
        \\
        %%%%%%%%%%%%%%%%%%
        \hline
        \textbf{Description}
         &
        Quản lý hàng hóa có thể thêm hàng vào kho
        \\
        %%%%%%%%%%%%%%%%%%
        \hline
        \textbf{Preconditions}
         &
        Quản lý hàng hóa đã đăng nhập
        \\
        %%%%%%%%%%%%%%%%%%
        \hline
        \textbf{Postconditions}
         &
        Hàng được thêm
        \\
        %%%%%%%%%%%%%%%%%%
        \hline
        \textbf{Trigger}
         &
        Quản lý hàng hóa chọn "Quản lý hàng hóa"
        \\
        %%%%%%%%%%%%%%%%%%
        \hline
        \begin{flushleft}
            \textbf{Normal flow}
        \end{flushleft}
         &
        1. Hệ thống hiển thị danh sách hàng hóa

        2. Quản lý hàng hóa tìm kiếm và chọn hàng muốn xuất

        3. Thông tin chi tiết hàng được hiển thị

        4. Quản lý hàng hóa chọn vào "Nhập kho"

        5. Quản lý hàng hóa nhập thông tin

        6. Quản lý hàng hóa chọn "Thực hiện"

        7. Hệ thống cập nhật thông tin
        \\
        %%%%%%%%%%%%%%%%%%
        \hline
        \begin{flushleft}
            \textbf{Alternative flow}
        \end{flushleft}
         &
        5.1. Quản lý hàng hóa chọn hủy
        \begin{quote}
            5.1.1 Quay lại bước 3
        \end{quote}
        \\
        %%%%%%%%%%%%%%%%%%
        \hline
        \textbf{Exceptions}
         &
        Không
        \\
        %%%%%%%%%%%%%%%%%%
        \hline
        \caption{Đặc tả use-case \textbf{Thêm hàng vào kho}}
    \end{longtable}
}


\subsubsection{Xuất hàng trong kho}
{
    \setlength\extrarowheight{6pt}
    \begin{longtable}{| p{.20\textwidth} | p{.80\textwidth} |}
        \hline
        \textbf{Use-case name}
         &
        Xuất hàng trong kho
        \\
        %%%%%%%%%%%%%%%%%%
        \hline
        \textbf{Actor}
         &
        Quản lý hàng hóa
        \\
        %%%%%%%%%%%%%%%%%%
        \hline
        \textbf{Description}
         &
        Quản lý hàng hóa có thể xuất hảng khỏi kho
        \\
        %%%%%%%%%%%%%%%%%%
        \hline
        \textbf{Preconditions}
         &
        Quản lý hàng hóa đã đăng nhập
        \\
        %%%%%%%%%%%%%%%%%%
        \hline
        \textbf{Postconditions}
         &
        Hàng được xuất khỏi kho
        \\
        %%%%%%%%%%%%%%%%%%
        \hline
        \textbf{Trigger}
         &
        Quản lý hàng hóa chọn "Quản lý hàng hóa"
        \\
        %%%%%%%%%%%%%%%%%%
        \hline
        \begin{flushleft}
            \textbf{Normal flow}
        \end{flushleft}
         &
        1. Hệ thống hiển thị danh sách hàng hóa

        2. Quản lý hàng hóa tìm kiếm và chọn hàng muốn xuất

        3.Thông tin chi tiết hàng được hiển thị

        4. Quản lý hàng hóa tìm kiếm và chọn kho muốn xuất

        5. Hệ thống hiển thị biểu mẫu thông tin quy trình

        6. Quản lý hàng hóa nhập thông tin

        7. Quản lý hàng hóa chọn "Thực hiện"

        8. Hệ thống cập nhật thông tin
        \\
        %%%%%%%%%%%%%%%%%%
        \hline
        \textbf{Alternative flow}
         &
        Không
        \\
        %%%%%%%%%%%%%%%%%%
        \hline
        \textbf{Exceptions}
         &
        Không
        \\
        %%%%%%%%%%%%%%%%%%
        \hline
        \caption{Đặc tả use-case \textbf{Xuất hàng trong kho}}
    \end{longtable}
}


\subsubsection{Chuyển kho}
{
    \setlength\extrarowheight{6pt}
    \begin{longtable}{| p{.20\textwidth} | p{.80\textwidth} |}
        \hline
        \textbf{Use-case name}
         &
        Chuyển kho
        \\
        %%%%%%%%%%%%%%%%%%
        \hline
        \textbf{Actor}
         &
        Quản lý hàng hóa
        \\
        %%%%%%%%%%%%%%%%%%
        \hline
        \textbf{Description}
         &
        Quản lý hàng hóa có thể chuyển hàng từ kho này đến kho khác
        \\
        %%%%%%%%%%%%%%%%%%
        \hline
        \textbf{Preconditions}
         &
        Quản lý hàng hóa đã đăng nhập
        \\
        %%%%%%%%%%%%%%%%%%
        \hline
        \textbf{Postconditions}
         &
        Hàng chuyển đi
        \\
        %%%%%%%%%%%%%%%%%%
        \hline
        \textbf{Trigger}
         &
        Quản lý hàng hóa chọn "Quản lý hàng hóa"
        \\
        %%%%%%%%%%%%%%%%%%
        \hline
        \begin{flushleft}
            \textbf{Normal flow}
        \end{flushleft}
         &
        1. Hệ thống hiển thị danh sách hàng hóa

        2. Quản lý hàng hóa tìm kiếm và chọn hàng muốn xuất

        3. Thông tin chi tiết hàng được hiển thị

        4. Quản lý hàng hóa tìm kiếm và chọn kho muốn xuất

        5. Hệ thống hiển thị biểu mẫu thông tin quy trình

        6. Quản lý hàng hóa nhập thông tin

        7. Quản lý hàng hóa chọn "Thực hiện"

        8. Hệ thống cập nhật thông tin
        \\
        %%%%%%%%%%%%%%%%%%
        \hline
        \textbf{Alternative flow}
         &
        Không
        \\
        %%%%%%%%%%%%%%%%%%
        \hline
        \textbf{Exceptions}
         &
        Không
        \\
        %%%%%%%%%%%%%%%%%%
        \hline
        \caption{Đặc tả use-case \textbf{Chuyển kho}}
    \end{longtable}
}

\subsubsection{Nhận trả hàng}
{
    \setlength\extrarowheight{6pt}
    \begin{longtable}{| p{.20\textwidth} | p{.80\textwidth} |}
        \hline
        \textbf{Use-case name}
         &
        Nhận trả hàng
        \\
        %%%%%%%%%%%%%%%%%%
        \hline
        \textbf{Actor}
         &
        Quản lý hàng hóa
        \\
        %%%%%%%%%%%%%%%%%%
        \hline
        \textbf{Description}
         &
        Quản lý hàng hóa có thể chuyển hàng từ kho này đến kho khác
        \\
        %%%%%%%%%%%%%%%%%%
        \hline
        \textbf{Preconditions}
         &
        Quản lý hàng hóa đã đăng nhập
        \\
        %%%%%%%%%%%%%%%%%%
        \hline
        \textbf{Postconditions}
         &
        Hàng chuyển đi
        \\
        %%%%%%%%%%%%%%%%%%
        \hline
        \textbf{Trigger}
         &
        Quản lý hàng hóa chọn "Quản lý hàng hóa"
        \\
        %%%%%%%%%%%%%%%%%%
        \hline
        \begin{flushleft}
            \textbf{Normal flow}
        \end{flushleft}
         &
        1. Hệ thống hiển thị danh sách hàng hóa

        2. Quản lý hàng hóa tìm kiếm và chọn hàng muốn xuất

        3. Thông tin chi tiết hàng được hiển thị

        4. Quản lý hàng hóa tìm kiếm và chọn kho muốn xuất

        5. Hệ thống hiển thị biểu mẫu thông tin quy trình

        6. Quản lý hàng hóa nhập thông tin

        7. Quản lý hàng hóa chọn "Thực hiện"

        8. Hệ thống cập nhật thông tin
        \\
        %%%%%%%%%%%%%%%%%%
        \hline
        \textbf{Alternative flow}
         &
        Không
        \\
        %%%%%%%%%%%%%%%%%%
        \hline
        \textbf{Exceptions}
         &
        Không
        \\
        %%%%%%%%%%%%%%%%%%
        \hline
        \caption{Đặc tả use-case \textbf{Nhận trả hàng}}
    \end{longtable}
}

\newpage

\subsection{Quản lý đơn hàng}
\begin{figure}[!htp]
    \centering
    \includegraphics[width=5in]{img/UseCase/UseCase-Quản lý đơn hàng.drawio.png}
    \newline
    \caption{Use case quản lý đơn hàng}
\end{figure}

\subsubsection{Xem thông tin đơn hàng trực tiếp}
{
    \setlength\extrarowheight{6pt}
    \begin{longtable}{| p{.20\textwidth} | p{.80\textwidth} |}
        \hline
        \textbf{Use-case name}
         &
        Xem thông tin đơn hàng trực tiếp
        \\
        %%%%%%%%%%%%%%%%%%
        \hline
        \textbf{Actor}
         &
        Quản lý chi nhánh, trưởng chi nhánh
        \\
        %%%%%%%%%%%%%%%%%%
        \hline
        \textbf{Description}
         &
        Quản lý chi nhánh, trưởng chi nhánh có thể xem thông tin các đơn hàng
        \\
        %%%%%%%%%%%%%%%%%%
        \hline
        \textbf{Preconditions}
         &
        Quản lý chi nhánh, trưởng chi nhánh đã đăng nhập
        \\
        %%%%%%%%%%%%%%%%%%
        \hline
        \textbf{Postconditions}
         &
        Nhận được thông tin đơn hàng
        \\
        %%%%%%%%%%%%%%%%%%
        \hline
        \textbf{Trigger}
         &
        Quản lý chi nhánh, trưởng chi nhánh chọn "Quản lý đơn hàng"
        \\
        %%%%%%%%%%%%%%%%%%
        \hline
        \begin{flushleft}
            \textbf{Normal flow}
        \end{flushleft}
         &
        1. Quản lý chi nhánh, trưởng chi nhánh chọn khoảng thời gian muốn xem

        2. Hệ thống hiển thị danh sách các đơn hàng

        3. Quản lý chi nhánh, trưởng chi nhánh tìm và chọn đơn muốn xem

        4. Hệ thống hiển thị thông tin đơn hàng
        \\
        %%%%%%%%%%%%%%%%%%
        \hline
        \textbf{Alternative flow}
         &
        Không
        \\
        %%%%%%%%%%%%%%%%%%
        \hline
        \textbf{Exceptions}
         &
        Không
        \\
        %%%%%%%%%%%%%%%%%%
        \hline
        \caption{Đặc tả use-case \textbf{Xem thông tin đơn hàng trực tiếp}}
    \end{longtable}
}

\subsubsection{Xem thông tin đơn hàng trực tuyến}
{
    \setlength\extrarowheight{6pt}
    \begin{longtable}{| p{.20\textwidth} | p{.80\textwidth} |}
        \hline
        \textbf{Use-case name}
         &
        Xem thông tin đơn hàng trực tuyến
        \\
        %%%%%%%%%%%%%%%%%%
        \hline
        \textbf{Actor}
         &
        Quản trị viên
        \\
        %%%%%%%%%%%%%%%%%%
        \hline
        \textbf{Description}
         &
        Quản trị viên có thể xem thông tin các đơn hàng trực tuyến
        \\
        %%%%%%%%%%%%%%%%%%
        \hline
        \textbf{Preconditions}
         &
        Quản trị viên đã đăng nhập
        \\
        %%%%%%%%%%%%%%%%%%
        \hline
        \textbf{Postconditions}
         &
        Nhận được thông tin đơn hàng
        \\
        %%%%%%%%%%%%%%%%%%
        \hline
        \textbf{Trigger}
         &
        Quản trị viên chọn "Đơn hàng trực tuyến" trong "Quản lý đơn hàng
        \\
        %%%%%%%%%%%%%%%%%%
        \hline
        \begin{flushleft}
            \textbf{Normal flow}
        \end{flushleft}
         &
        1. Quản trị viên chọn khoảng thời gian muốn xem

        2. Hệ thống hiển thị danh sách các đơn hàng

        3. Quản trị viên tìm và chọn đơn muốn xem

        4. Hệ thống hiển thị thông tin đơn hàng
        \\
        %%%%%%%%%%%%%%%%%%
        \hline
        \textbf{Alternative flow}
         &
        Không
        \\
        %%%%%%%%%%%%%%%%%%
        \hline
        \textbf{Exceptions}
         &
        Không
        \\
        %%%%%%%%%%%%%%%%%%
        \hline
        \caption{Đặc tả use-case \textbf{Xem thông tin đơn hàng trực tuyến}}
    \end{longtable}
}

\subsubsection{Cập nhật trạng thái đơn hàng}
{
    \setlength\extrarowheight{6pt}
    \begin{longtable}{| p{.20\textwidth} | p{.80\textwidth} |}
        \hline
        \textbf{Use-case name}
         &
        Cập nhật trạng thái đơn hàng.
        \\
        \hline
        \textbf{Actor}
         &
        Giao vận.
        \\
        \hline
        \textbf{Description}
         &
        Bên giao vận gửi trạng thái cập nhật đơn hàng khi có sự thay đổi trạng thái ngoài thực tế.
        \\
        \hline
        \textbf{Preconditions}
         &
        Bên giao vận gửi trạng thái cập nhật của đơn hàng.
        \\
        \hline
        \textbf{Postconditions}
         &
        Hệ thống ghi nhận trạng thái mới của đơn hàng.
        \\
        \hline
        \textbf{Trigger}
         &
        Không.
        \\
        \hline
        \textbf{Normal flow}
         &
        1. Bên giao vận gửi trạng thái cập nhật của đơn hàng

        2. Hệ thống ghi nhận cập nhật lại trạng thái mới cho đơn hàng.
        \\
        \hline
        \textbf{Alternative flow}
         &
        Không
        %%%%%%%%%%%%%%%%
        \\
        \hline
        \textbf{Exceptions}
         &
        Không
        %%%%%%%%%%%%%%%%%%%%
        \\
        \hline
        \caption{Đặc tả use-case \textbf{cập nhật trạng thái đơn hàng}}
    \end{longtable}
}

\subsubsection{Xem trạng thái đơn hàng}
{
    \setlength\extrarowheight{6pt}
    \begin{longtable}{| p{.20\textwidth} | p{.80\textwidth} |}
        \hline
        \textbf{Use-case name}
         &
        Xem trạng thái đơn hàng.
        \\
        \hline
        \textbf{Actor}
         &
        Khách hàng.
        \\
        \hline
        \textbf{Description}
         &
        Khách hàng xem trạng thái hiện tại của đơn hàng.
        \\
        \hline
        \textbf{Preconditions}
         &
        Khách hàng đã đăng nhập.
        \\
        \hline
        \textbf{Postconditions}
         &
        Khách hàng xem được trạng thái hiện tại của đơn hàng mà khách hàng muốn xem.
        \\
        \hline
        \textbf{Trigger}
         &
        Không.
        \\
        \hline
        \begin{flushleft}
            \textbf{Normal flow}

        \end{flushleft}
         &
        1. Khách hàng chọn vào mục "Quản lý đơn hàng".

        2. Hệ thống hiển thị danh sách các đơn hàng ở tab mặc định chưa giao hàng.

        3. Khách hàng lựa chọn đơn hàng mà mình muốn xem thông tin và nhấn vào đơn hàng đó.

        4. Hệ thống hiển thị thông tin chi tiết về trạng thái đơn hàng.

        5. Người dùng xem kết quả trạng thái đơn hàng mà mình muốn.
        \\
        \hline
        \begin{flushleft}
            \textbf{Alternative flow}
        \end{flushleft}
         &
        3.1. Khách hàng muốn xem các đơn hàng đã giao.
        \begin{quote}
            3.1.1. Khách hàng chọn mục "Đơn đã giao".

            3.1.2. Hệ thống hiển thị danh sách các đơn hàng đã giao của khách hàng.

            3.1.3. Quay lại bước 3.
        \end{quote}
        %%%%%%%%%%%%%%%%
        \\
        \hline
        \textbf{Exceptions}
         &
        Không
        %%%%%%%%%%%%%%%%%%%%
        \\
        \hline
        \caption{Đặc tả use-case \textbf{xem trạng thái đơn hàng}}
    \end{longtable}
}


\subsubsection{Yêu cầu xuất hóa đơn}
{
    \setlength\extrarowheight{6pt}
    \begin{longtable}{| p{.20\textwidth} | p{.80\textwidth} |}
        \hline
        \textbf{Use-case name}
         &
        Yêu cầu xuất hóa đơn.
        \\
        \hline
        \textbf{Actor}
         &
        Khách hàng.
        \\
        \hline
        \textbf{Description}
         &
        Khách hàng yêu cầu hệ thống xuất hóa đơn cho đơn hàng online.
        \\
        \hline
        \textbf{Preconditions}
         &
        Khách hàng đã từng mua hàng thành công trên trang bán hàng online.
        \\
        \hline
        \textbf{Postconditions}
         &
        Hệ thống gửi khách hàng hóa đơn của đơn hàng tương ứng.
        \\
        \hline
        \textbf{Trigger}
         &
        Không.
        \\
        \hline
        \begin{flushleft}
            \textbf{Normal flow}
        \end{flushleft}
         &
        1. Khách hàng chọn vào mục "Quản lý đơn hàng".

        2. Hệ thống hiển thị danh sách các đơn hàng ở tab mặc định chưa giao hàng.

        3. Khách hàng chọn mục "Đơn đã giao".

        4. Hệ thống hiển thị danh sách các đơn hàng đã giao của khách hàng.

        5. Khách hàng lựa chọn đơn hàng mà mình muốn xem thông tin và nhấn vào đơn hàng đó sau đó chọn "Yêu cầu xuất hóa đơn".

        6. Hệ thống thông báo kết quả và trả về tệp hóa đơn cho khách hàng.

        7. Người dùng xem kết quả trong tệp hóa đơn.
        \\
        \hline
        \textbf{Alternative flow}
         &
        Không
        %%%%%%%%%%%%%%%%
        \\
        \hline
        \begin{flushleft}
            \textbf{Exceptions}
        \end{flushleft}
         &
        3.1. Hệ thống tạo hóa đơn thất bại.
        \begin{quote}
            3.1.1. Quay lại bước 3.
        \end{quote}
        %%%%%%%%%%%%%%%%%%%%
        \\
        \hline
        \caption{Đặc tả use-case \textbf{yêu cầu xuất hóa đơn}}
    \end{longtable}
}




\newpage
% BPMN
% % !TeX root = ..\main.tex
\section{Thiết kế lược đồ BPMN}

Các lược đồ BPMN được nhóm thiết kế trên app Bizagi.
\subsection{Khách hàng mua hàng}
    \begin{figure}[!htp]
        \centering
        \includegraphics[width=17cm]{img/BPMN/customer_buy/customer_buy.png}
        \newline
        \caption{Lược đồ BPMN cho quy trình khách hàng mua hàng}
    \end{figure}
    \textbf{Mô tả:}
    \begin{itemize}
        \item \textbf{Tìm sản phẩm}: Đây là 1 quy trình con chứa quy trình tìm kiếm sản phẩm của  khách hàng.
        \item \textbf{Xem xét thông tin chi tiết}: Đây là 1 quy trình con chứa quy trình xem thông tin chi tiết sản phẩm của  khách hàng.
        \item \textbf{Thêm sản phẩm vào giỏ hàng}: Đây là 1 quy trình con chứa quy trình thêm sản phẩm vào giỏ hàng.
        \item \textbf{Thêm nhanh vào giỏ hàng}: Đây là 1 quy trình con chứa quy trình thêm nhanh sản phẩm vào giỏ hàng.
        \item \textbf{Quản lý giỏ hàng}: Đây là 1 quy trình con chứa quy trình quản lý giỏ hàng của khách hàng.
        \item \textbf{Mua hàng}: Đây là 1 quy trình con chứa quy trình của việc mua hàng trong các giỏ hàng khách hàng.
    \end{itemize}
    \begin{figure}[!htp]
        \centering
        \includegraphics[width=13cm]{img/BPMN/customer_buy/customer_search_product.png}
        \newline
        \caption{Lược đồ BPMN cho quy trình con tìm kiếm sản phẩm}
    \end{figure}
    \textbf{Mô tả:}
    \begin{itemize}
        \item \textbf{Chọn loại sản phầm}: Đây là 1 quy trình con chứa quy trình chọn loại sản phẩm mà người dùng muốn tình kiếm.
    \end{itemize}
    \begin{figure}[!htp]
        \centering
        \includegraphics[width=12cm]{img/BPMN/customer_buy/customer_select_type.png}
        \newline
        \caption{Lược đồ BPMN cho quy trình con chọn loại sản phẩm mà khác hàng muốn tìm}
    \end{figure}
    \begin{figure}[!htp]
        \centering
        \includegraphics[width=12cm]{img/BPMN/customer_buy/customer_product_detail.png}
        \newline
        \caption{Lược đồ BPMN cho quy trình xem xét thông tin chi tiết của sản phẩm}
    \end{figure}
    \begin{figure}[!htp]
        \centering
        \includegraphics[width=15cm]{img/BPMN/customer_buy/customer_add_to_card.png}
        \newline
        \caption{Lược đồ BPMN cho quy trình thêm sản phẩm vào giỏ hàng}
    \end{figure}  
    \begin{figure}[!htp]
        \centering
        \includegraphics[width=15cm]{img/BPMN/customer_buy/customer_add_fast.png}
        \newline
        \caption{Lược đồ BPMN cho quy trình thêm nhanh sản phẩm vào giỏ hàng}
    \end{figure}
    \begin{figure}[!htp]
        \centering
        \includegraphics[width=15cm]{img/BPMN/customer_buy/customer_cart.png}
        \newline
        \caption{Lược đồ BPMN cho quy trình quản lý giỏ hàng}
    \end{figure}
    \begin{figure}[!htp]
        \centering
        \includegraphics[width=15cm]{img/BPMN/customer_buy/customer_buy_order.png}
        \newline
        \caption{Lược đồ BPMN cho quy trình mua hàng}
    \end{figure}
    \newpage
    \textbf{Mô tả:}
    \begin{itemize}
        \item \textbf{Xử lý giao hàng và giá tiền}: Đây là 1 quy trình con chứa quy trình con xử lý về dịch vụ giao hàng cho đơn hàng và tính toán tất cả các chi phí.
        \item \textbf{Thanh toán}: Đây là 1 quy trình con chứa quy trình thanh toán đơn hàng của khách hàng.
    \end{itemize}
    \begin{figure}[!htp]
        \centering
        \includegraphics[width=5in]{img/BPMN/customer_buy/customer_calc_fee.png}
        \newline
        \caption{Lược đồ BPMN cho quy trình tính toán chi phí đơn hàng}
    \end{figure}
    \begin{figure}[!htp]
        \centering
        \includegraphics[width=15cm]{img/BPMN/customer_buy/customer_payment.png}
        \newline
        \caption{Lược đồ BPMN cho quy trình thanh toán đơn hàng}
    \end{figure}

\subsection{Quản lý sự kiện}

\subsubsection{Thêm sự kiện}

\begin{figure}[!htp]
    \centering
    \includegraphics[width=14cm]{img/BPMN/event/add_event.png}
    \newline
    \caption{Lược đồ BPMN cho quy trình thêm sự kiện}
\end{figure}

\subsubsection{Sửa sự kiện}

\begin{figure}[!htp]
    \centering
    \includegraphics[width=14cm]{img/BPMN/event/edit_event.png}
    \newline
    \caption{Lược đồ BPMN cho quy trình chỉnh sửa sự kiện}
\end{figure}

\subsubsection{Xóa sự kiện}

\begin{figure}[!htp]
    \centering
    \includegraphics[width=14cm]{img/BPMN/event/delete_event.png}
    \newline
    \caption{Lược đồ BPMN cho quy trình xóa sự kiện}
\end{figure}

\subsection{Khách hàng mua hàng trực tiếp}
\begin{figure}[!htp]
    \centering
    \includegraphics[width=14cm]{img/BPMN/Hien/Customer_buyOffline.png}
    \newline
    \caption{Lược đồ BPMN cho quy trình khách hàng mua hàng trực tiếp}
\end{figure}

\subsection{Đăng nhập}
\begin{figure}[!htp]
    \centering
    \includegraphics[width=14cm]{img/BPMN/Hien/Customer_login.png}
    \newline
    \caption{Lược đồ BPMN cho quy trình đăng nhập}
\end{figure}

\subsection{Đăng ký tài khoản}
\begin{figure}[!htp]
    \centering
    \includegraphics[width=14cm]{img/BPMN/Hien/Customer_register.png}
    \newline
    \caption{Lược đồ BPMN cho quy trình đăng ký tài khoản mới}
\end{figure}


\subsection{Làm mới mật khẩu}
\begin{figure}[!htp]
    \centering
    \includegraphics[width=14cm]{img/BPMN/Hien/Customer_resetPassword.png}
    \newline
    \caption{Lược đồ BPMN cho quy trình làm mới mật khẩu}
\end{figure}



\subsection{Quản lý chi nhánh}
\begin{figure}[!htp]
    \centering
    \includegraphics[width=14cm]{img/BPMN/Hien/Branch_management.png}
    \newline
    \caption{Lược đồ BPMN cho quy trình quản lý chi nhánh}
\end{figure}


\subsection{Quản lý nhân viên}
\begin{figure}[!htp]
    \centering
    \includegraphics[width=14cm]{img/BPMN/Hien/Employee_Management.png}
    \newline
    \caption{Lược đồ BPMN cho quy trình quản lý nhân viên}
\end{figure}

\subsubsection*{Tạo yêu cầu thêm, xóa nhân viên}
\begin{figure}[!htp]
    \centering
    \includegraphics[width=14cm]{img/BPMN/Hien/Employee_request.png}
    \newline
    \caption{Lược đồ BPMN cho quy trình quản lý nhân viên}
\end{figure}



\newpage
% Architecture
% % !TeX root = ..\main.tex
\section{Kiến trúc hệ thống}
\subsection{Tổng quan}
\begin{figure}[!htp]
	\centering
	\includegraphics[width=6in]{img/Architecture/general-architect.png}
	\newline
	\caption{Tổng quan kiến trúc hệ thống}
\end{figure}

\subsection{Tầng UI}

\subsubsection{CustomerUI}
Module này bao gồm các class dùng để hiển thị UI đối với phía khách hàng
\begin{figure}[!htp]
	\centering
	\includegraphics[width=17cm]{img/Architecture/customer UI.png}
	\newline
	\caption{Lược đồ class của Module CustomerUI}
\end{figure}
\textbf{Mô tả:}
\begin{quote}
	\begin{itemize}
		\item CustomerUI
		\item NavBar
		\item Header
		\item Footer
		\item ..
	\end{itemize}
\end{quote}

\subsubsection{LoginUI}
LoginUI là một class dùng để hiển thị trang đăng nhập cho cả phía khách hàng và quản trị viên
\begin{figure}[!htp]
	\centering
	\includegraphics[width=2cm]{img/Architecture/loginUI.png}
	\newline
	\caption{Lược đồ class của LoginUI}
\end{figure}
\textbf{Mô tả:}
\begin{quote}
	\begin{itemize}
		\item CustomerUI
		\item NavBar
		\item Header
		\item Footer
		\item ..
	\end{itemize}
\end{quote}

\subsubsection{AdminUI}
Module này bao gồm các class dùng để hiển thị UI đối với phía quản trị viên

\begin{figure}[!htp]
	\centering
	\includegraphics[width=17cm]{img/Architecture/adminUI.png}
	\newline
	\caption{Lược đồ class của Module AdminUI}
\end{figure}
\textbf{Mô tả:}
\begin{quote}
	\begin{itemize}
		\item CustomerUI
		\item NavBar
		\item Header
		\item Footer
		\item ..
	\end{itemize}
\end{quote}

\subsection{Tầng API}
\subsubsection{Customer APIs}
Module này dùng để chứa các API được xuất ra ngoài cho phía khách hàng

\subsubsection{Admin APIs}
Module này dùng để chứa các API được xuất ra ngoài cho phía quản trị viên

\subsubsection{Shared APIs}
Module này dùng để chứa các API được xuất ra ngoài cho bất kì đối tượng người dùng nào


\subsection{Tầng service}

\subsubsection{}


\newpage
% Database
% % !TeX root = ..\main.tex
\section{Thiết kế cơ sở dữ liệu}
\par Nhóm lựa chọn kiến trúc Microservice cho hệ thống, với việc chia ra các service nhỏ và mỗi service có một database riêng phục vụ cho nghiệp vụ của riêng service đó, nên ở đây các database được thiết kế riêng lẻ cho từng service.

\subsection{ERD/EERD}
\subsubsection{Service Khách hàng}
\begin{figure}[!htp]
    \begin{center}
        \includegraphics[width=1\textwidth]{img/database/erd/eerd-customer.png}
        \newline
        \caption{ERD cho database của service Khách hàng}
    \end{center}
\end{figure}

\par Database của service Khách hàng bao gồm hai thực thể là Khách hàng và Tài khoản của khách hàng. Thực thể khách hàng chứa các thông tin cá nhân của khách hàng, bao gồm các thuộc tính:
\begin{itemize}
    \item Mã khách hàng: thuộc tính khóa, là mã của khách hàng.
    \item Tên: Là tên của khách hàng.
    \item Giới tính: Là giới tính của khách hàng.
    \item Email: Là email của khách hàng.
    \item Số điện thoại: Là số điện thoại của khách hàng.
    \item Số nhà, đường: Là địa chỉ sô nhà, đường của khách hàng ở.
    \item Phường / xã: Là địa chỉ phường / xã của khách hàng ở.
    \item Quận / huyện: Là địa chỉ quận / huyện của khách hàng ở.
    \item Tỉnh / thành phố: Là địa chỉ tỉnh / thành phố của khách hàng ở.
\end{itemize}

\par Thực thể Tài khoản của khách hàng bao gồm các thuộc tính:
\begin{itemize}
    \item Tên đăng nhập: tên đăng nhập của tài khoản của khách hàng.
\end{itemize}

\par Khách hàng và Tài khoản có mối quan hệ thuộc, đây là mối quan hệ bắt buộc, 1-1: một khách hàng chỉ thuộc 1 tài khoản và 1 tài khoản chỉ dùng cho 1 khách hàng.

\subsubsection{Service Giỏ hàng}
\begin{figure}[!htp]
    \begin{center}
        \includegraphics[width=1\textwidth]{img/database/erd/eerd-cart.png}
        \newline
        \caption{ERD cho database của service Giỏ hàng}
    \end{center}
\end{figure}

\par Database của service Giỏ hàng bao gồm hai thực thể là Giỏ hàng và Hàng trong giỏ hàng. Thực thể Giỏ hàng chứa thông tin về giỏ hàng, bao gồm các thuộc tính:
\begin{itemize}
    \item Mã giỏ hàng: thuộc tính khóa, là mã của giỏ hàng.
    \item Mã khách hàng: là mã khách hàng của khách hàng chủ của giỏ hàng đó.
    \item Số lượng món: là thuộc tính dẫn xuất, số lượng loại sản phẩm có trong giỏ hàng.
\end{itemize}

\par Thực thể Hàng trong giỏ hàng chứa thông tin của mặt hàng trong trong giỏ hàng đó. Thực thể này bao gồm các thuộc tính:
\begin{itemize}
    \item Mã hàng: thuộc tính khóa, là mã của loại hàng đó.
    \item Kích thước: thuộc tính khóa, là kích thước của loại hàng đó.
    \item Màu: thuộc tính khóa, là màu sắc của loại hàng đó.
\end{itemize}

\par Giỏ hàng và Hàng có mối quan hệ chứa, Mối quan hệ là M-N. Giỏ hàng tham gia vào mối quan hệ với ràng buộc tham gia không bắt buộc: Giỏ hàng có thể chứa hàng hoặc không, có thể chứa hàng với số lượng khác nhau. Hàng tham gia vào mối quan hệ với ràng buộc tham gia không bắt buộc, một hàng có thể có hoặc không thuộc giỏ hàng nào. Mối quan hệ này có các thuộc tính:
\begin{itemize}
    \item Số lượng: số lượng của món hàng có trong giỏ hàng.
\end{itemize}

\subsubsection{Service Đơn hàng của khách hàng}
\begin{figure}[!htp]
    \begin{center}
        \includegraphics[width=1\textwidth]{img/database/erd/eerd-customer-order.png}
        \newline
        \caption{ERD cho database của service Đơn hàng của khách hàng}
    \end{center}
\end{figure}

\par Database của service Đơn hàng của khách hàng bao gồm hai thực thể là Đơn hàng và Hàng trong đơn hàng. Thực thể Đơn hàng chứa thông tin về đơn hàng, bao gồm các thuộc tính:
\begin{itemize}
    \item Mã đơn hàng: thuộc tính khóa, là mã của đơn hàng.
    \item Mã khách hàng: là mã khách hàng của khách hàng chủ của đơn hàng đó.
    \item Ngày giao dịch: là ngày thực hiện giao dịch của đơn hàng.
    \item Phí ship: là phí vận chuyển của đơn hàng.
    \item Mong đợi: là thời gian ước tính để giao đơn hàng.
    \item Tổng giá tiền: là thuộc tính dẫn xuất, là tổng giá tiền của đơn hàng.
    \item Thông tin trạng thái: là thuộc tính đa trị, là trạng thái của đơn hàng, bao gồm 2 thuộc tính con là Trạng thái và Thời gian. Trạng thái là trạng thái của đơn hàng và thời gian là thời điểm xảy ra trạng thái của đơn hàng.
\end{itemize}

\par Thực thể Hàng trong đơn hàng chứa thông tin của mặt hàng trong trong đơn hàng đó. Thực thể này bao gồm các thuộc tính:
\begin{itemize}
    \item Mã hàng: thuộc tính khóa, là mã của loại hàng đó.
    \item Kích thước: thuộc tính khóa, là kích thước của loại hàng đó.
    \item Màu: thuộc tính khóa, là màu sắc của loại hàng đó.
    \item Số lượng: là số lượng của loại hàng trong đơn hàng.
    \item Đơn giá: là giá tiền của loại hàng trong đơn hàng.
    \item Thuế: là thuế được thu trên loại hàng đó trong đơn hàng.
    \item Thành tiền: là thuộc tính dẫn xuất, giá tiền cuối cùng được tính cho loại hàng đó trong đơn hàng.
\end{itemize}

\par Đơn hàng và Hàng có mối quan hệ Bao gồm, đây là quan hệ bắt buộc, Mối quan hệ là 1-N: đơn hàng bắt buộc có hàng, và có thể có nhiều hàng; hàng trong đơn hàng bắt buộc phải thuộc về một đơn hàng, và chỉ thuộc về duy nhất một đơn hàng.

\subsubsection{Service Sự kiện}
\begin{figure}[!htp]
    \begin{center}
        \includegraphics[width=1\textwidth]{img/database/erd/eerd-event.png}
        \newline
        \caption{ERD cho database của service Sự kiện}
    \end{center}
\end{figure}

\par Database của service Sự kiện bao gồm hai thực thể là Sự kiện và Hàng. Thực thể Sự kiện chứa thông tin về sự kiện, bao gồm các thuộc tính:
\begin{itemize}
    \item Mã sự kiện: thuộc tính khóa, là mã của sự kiện.
    \item Tên: là của sự kiện.
    \item Giảm giá: là giá trị giảm giá của sự kiện đối với các mặt hàng trong sự kiện.
    \item Thời gian bắt đầu: là thời gian bắt đầu sự kiện.
    \item Thời gian kết thúc: là thời gian kết thúc sự kiện.
\end{itemize}

\par Thực thể Hàng trong sự kiện chứa thông tin của mặt hàng trong trong sự kiện đó. Thực thể này bao gồm các thuộc tính:
\begin{itemize}
    \item Mã hàng: thuộc tính khóa, là mã của loại hàng đó.
\end{itemize}

\par Sự kiện và Hàng có mối quan hệ có, mối quan hệ là M-N: sự kiện có thể có nhiều hàng, ít hàng hoặc không có hàng; hàng trong cũng có thể thuộc không hoặc ít hoặc nhiều sự kiện.

\subsubsection{Service Hàng hóa}
\begin{figure}[!htp]
    \begin{center}
        \includegraphics[width=1\textwidth]{img/database/erd/shop_online-hàng.png}
        \newline
        \caption{ERD cho database của service Hàng hóa}
    \end{center}
\end{figure}

\par Database của service Hàng hóa bao gồm hai thực thể là Hàng hóa và Hàng trong kho. Thực thể hàng hóa chỉ thông tin các mặt hàng được bán, bao gồm các thuộc tính:
\begin{itemize}
    \item Mã hàng: thuộc tính khóa, là mã của mặt hàng.
    \item Kích thước: thuộc tính khóa, là kích thước của mặt hàng.
    \item Màu: thuộc tính khóa, là màu của mặt hàng.
    \item Tên hàng: là tên của mặt hàng.
    \item Lứa tuổi: là lứa tuổi phù hợp của mặt hàng.
    \item Giới tính: là giới tính phù hợp của mặt hàng.
    \item Nhà sản xuất: là nhà sản xuất của mặt hàng.
    \item Hình ảnh: thuộc tính đa trị, bao gồm các hình ảnh của mặt hàng.
    \item Được bán: cho biết mặt hàng có đang được mở bán hay không.
    \item Đơn giá: giá tiền cho một đơn vị của mặt hàng.
    \item Mô tả: thông tin thêm về mặt hàng.
    \item Loại: thể loại của mặt hàng.
    \item Số lượng hiện tại: thuộc tính dẫn xuất, là số lượng còn tồn hiện tại của mặt hàng.
\end{itemize}

\par Thực thể Hàng trong kho chỉ thông tin của mặt hàng trong một kho nào đó. Thực thể này bao gồm các thuộc tính:
\begin{itemize}
    \item Mã kho: thuộc tính khóa, là mã của kho chứa hàng đó.
    \item Số lượng: số lượng hàng đó còn trong kho.
    \item Ngày tạo: ngày mặt hàng lần đầu vào kho.
    \item Ngày thay đổi: ngày hàng trong kho có sự thay đổi.
\end{itemize}

\par Hàng hóa và Hàng trong kho có mối quan hệ Bao gồm, đây là mối quan hệ bắt buộc, 1-N: hàng hóa bắt buộc phải nằm trong kho, và một mặt hàng có thể nằm ở nhiều kho; hàng trong kho bắt buộc thuộc về một mặt hàng, và chỉ thuộc về duy nhất một mặt hàng.

\subsubsection{Service Đơn hàng}
\begin{figure}[!htp]
    \begin{center}
        \includegraphics[width=1\textwidth]{img/database/erd/shop_online-đơnhàng.png}
        \newline
        \caption{ERD cho database của service Đơn hàng}
    \end{center}
\end{figure}

\par Database của service Đơn hàng bao gồm các thực thể Đơn hàng và Hàng hóa, chỉ các hàng hóa trong đơn hàng. Đơn hàng được cụ thể hóa xuống Đơn hàng online và Đơn hàng cửa hàng.
\par Thực thể Đơn hàng gồm các thuộc tính:
\begin{itemize}
    \item Mã đơn hàng: thuộc tính khóa, chỉ mã của đơn hàng.
    \item Ngày giao dịch: ngày thực hiện giao dịch mua hàng và sinh ra đơn hàng.
    \item Tổng giá tiền: thuộc tính dẫn xuất, là tổng giá tiền của đơn hàng.
\end{itemize}

\par Thực thể Đơn hàng online có thêm các thuộc tính:
\begin{itemize}
    \item Mã khách hàng: là mã khách hàng đã mua hàng online.
    \item Phí ship: giá tiền cho việc vận chuyển đơn hàng.
    \item Mong đợi: ngày giao hàng mong đợi.
    \item Thông tin trạng thái: thuộc tính phức hợp, bao gồm các thông tin Trạng thái và Thời gian chỉ các trạng thái và thời gian đạt đến trạng thái đó của đơn hàng.
\end{itemize}

\par Thực thể Đơn hàng cửa hàng có thêm thuộc tính:
\begin{itemize}
    \item Mã cửa hàng: mã cửa hàng bán hàng trong đơn hàng.
\end{itemize}

\par Thực thể Hàng bao gồm các thuộc tính:
\begin{itemize}
    \item Mã hàng: thuộc tính khóa, là mã của mặt hàng trong đơn hàng.
    \item Kích thước: thuộc tính khóa, là kích thước của hàng trong đơn hàng.
    \item Màu: thuộc tính khóa, là màu của hàng trong đơn hàng.
    \item Số lượng: số lượng hàng đã mua trong đơn hàng.
    \item Đơn giá: giá tiền một đơn vị của hàng trong đơn hàng.
    \item Thuế: giá tiền thuế của hàng trong đơn hàng.
    \item Thành tiền: thuộc tính dẫn xuất, là tổng giá tiền (đơn giá * số lượng) + thuế.
\end{itemize}

\par Đơn hàng và Hàng có quan hệ Bao gồm, đây là quan hệ bắt buộc, 1-N: đơn hàng bắt buộc có hàng, và có thể có nhiều hàng; hàng trong đơn hàng bắt buộc phải thuộc về một mặt hàng, và chỉ thuộc về duy nhất một mặt hàng.

\subsubsection{Service Kho}
\begin{figure}[!htp]
    \begin{center}
        \includegraphics[width=1\textwidth]{img/database/erd/shop_online-Kho.png}
        \newline
        \caption{ERD cho database của service Kho}
    \end{center}
\end{figure}

\par Database của service Kho bao gồm thực thể Kho và thực thể Nhân viên kho. Thực thể Kho bao gồm các thuộc tính:
\begin{itemize}
    \item Mã kho: thuộc tính khóa, là mã của kho.
    \item Tên kho: tên của kho.
    \item Ngày thành lập: ngày mà kho đi vào sử dụng.
    \item Địa chỉ: thuộc tính kết hợp, bao gồm các thuộc tính Số nhà, đường, Phường, Thành phố / Quận / Huyện, Tỉnh / Thành phố.
    \item Sức chứa: sức chứa hàng tối đa của kho.
    \item Còn lại: thuộc tính dẫn xuất, chỉ sức chứa còn lại mà kho có thể chứa hàng thêm.
\end{itemize}

\par Thực thể Nhân viên kho bao gồm các thuộc tính:
\begin{itemize}
    \item Mã nhân viên: thuộc tính khóa, là mã của nhân viên làm việc ở kho.
\end{itemize}

\par Kho và nhân viên kho có mối quan hệ Bao gồm là mối quan hệ bắt buộc, M-N: kho bắt buộc có nhân viên, và có thể có nhiều nhân viên làm việc; nhân viên buộc phải làm việc ở ít nhất một kho, và có thể làm việc ở nhiều kho. Mối quan hệ này có các thuộc tính:
\begin{itemize}
    \item Ngày bắt đầu: ngày nhân viên bắt đầu làm việc tại kho đó.
    \item Ngày kết thúc: ngày nhân viên nghỉ làm việc tại kho đó.
\end{itemize}

\par Mỗi kho có một quản lý, nên kho và nhân viên quản lý có mối quan hệ Được quản lý bởi. Mối quan hệ là M-N. Kho tham gia vào mối quan hệ với ràng buộc tham gia bắt buộc: kho phải có quản lý, và theo thời gian có nhiều quản lý theo từng thời gian khác nhau. Quản lý - cũng là nhân viên kho, tham gia vào mối quan hệ với ràng buộc tham gia không bắt buộc, một nhân viên không cần phải làm quản lý, và qua thời gian có thể làm quản lý ở nhiều kho. Mối quan hệ này có các thuộc tính:
\begin{itemize}
    \item Ngày bắt đầu: ngày quản lý bắt đầu làm việc tại kho đó.
    \item Ngày kết thúc: ngày quản lý nghỉ làm việc tại kho đó.
\end{itemize}

\subsection{Ánh xạ}
\par Từ các ERD/EERD trên, chúng ta có các ánh xạ sau tương ứng với mỗi service.

\subsubsection{Service Hàng hóa}
\begin{figure}[!htp]
    \begin{center}
        \includegraphics[width=1\textwidth]{img/database/mapping/mapping-hàng.png}
        \newline
        \caption{Ánh xạ cho database của service Hàng hóa}
    \end{center}
\end{figure}

\subsubsection{Service Đơn hàng}
\begin{figure}[!htp]
    \begin{center}
        \includegraphics[width=1\textwidth]{img/database/mapping/mapping-đơnhàng.png}
        \newline
        \caption{Ánh xạ cho database của service Đơn hàng}
    \end{center}
\end{figure}

\subsubsection{Service Kho}
\begin{figure}[!htp]
    \begin{center}
        \includegraphics[width=1\textwidth]{img/database/mapping/mapping-kho.png}
        \newline
        \caption{Ánh xạ cho database của service Kho}
    \end{center}
\end{figure}

\newpage
% UI
% % !TeX root = ..\main.tex
\section{Thiết kế giao diện}

\subsection{Giao diện chung}
\subsubsection{Đăng nhập}
\begin{figure}[!htp]
    \centering
    \includegraphics[width=5in]{img/UI/customer/login.png}
    \label{1}
    \newline
    \caption{Giao diện đăng nhập}
\end{figure}
\textbf{Mô tả:}
\begin{quote}
    \begin{enumerate}
        \item Chọn để chuyển tới trang đăng ký tài khoản.
        \item Nhập tên đăng nhập của tài khoản, yêu cầu tối thiểu 1 ký tự.
        \item Nhập mật khẩu của tài khoản, yêu cầu mật khẩu tối thiểu 6 ký tự.
        \item Chọn để thực hiện đăng nhập tài khoản, nếu thông tin tài khoản đúng thì sẽ được đưa đến giao diện mặc định cho tài khoản.
        \item Chọn "Quên mật khẩu" khi không nhớ mật khẩu của tài khoản để chuyển trang sang trang lấy lại mật khẩu.
    \end{enumerate}
\end{quote}


\subsection{Giao diện người dùng}
\subsubsection{Đăng ký}
\begin{figure}[!htp]
    \centering
    \includegraphics[width=5in]{img/UI/customer/register.png}
    \label{2}
    \newline
    \caption{Giao diện đăng ký tài khoản người dùng}
\end{figure}
\textbf{Mô tả:}
\begin{quote}
    \begin{enumerate}
        \item Chọn để chuyển tới trang đăng đăng nhập.
        \item Nhập tên đăng nhập của tài khoản, yêu cầu tối thiểu một ký tự.
        \item Nhập tên của người dùng, yêu cầu tối thiểu một ký tự.
        \item Nhập số điện thoại của người dùng, yêu cầu đúng định dạng số điện thoại 10 chữ số, bắt đầu bằng số 09 / 03 / 05 / 07 / 08.
        \item Nhập mật khẩu của tài khoản, yêu cầu tối thiểu 6 ký tự.
        \item Nhập lại mật khẩu của tài khoản, yêu cầu phải giống mật khẩu của tài khoản đã nhập.
        \item Chọn để thực hiện đăng ký tài khoản, nếu đăng ký thành công thì chuyển tới trang đăng nhập.
    \end{enumerate}
\end{quote}


\subsubsection{Tạo lại mật khẩu}
\begin{figure}[!htp]
    \centering
    \includegraphics[width=5in]{img/UI/customer/reset_new_password.png}
    \label{3}
    \newline
    \caption{Giao diện tạo lại mật khẩu người dùng}
\end{figure}
\textbf{Mô tả:}
\begin{quote}
    \begin{enumerate}
        \item Người dùng nhập số điện thoại của tài khoản, yêu cầu đúng định dạng số điện thoại 10 chữ số, bắt đầu bằng số 09 / 03 / 05 / 07 / 08.
        \item Người dùng chọn để hệ thống gửi xác nhận tin nhắn OTP và chuyển tới trang xác nhận OTP.
    \end{enumerate}
\end{quote}
\newpage
\begin{figure}[!htp]
    \centering
    \includegraphics[width=5in]{img/UI/customer/confirm_otp.png}
    \label{4}
    \newline
    \caption{Giao diện xác nhận OTP}
\end{figure}
\textbf{Mô tả:}
\begin{quote}
    \begin{enumerate}
        \item Người dùng nhập OTP đã được gửi qua tin nhắn, yêu cầu 6 ký tự.
        \item Người dùng chọn để hệ thống xác nhận và thông tin hợp lệ sẽ được chuyển đến trang tạo lại mật khẩu mới.
    \end{enumerate}
\end{quote}
\textbf{Mô tả:}
\begin{quote}
    \begin{enumerate}
        \item Người dùng nhập số điện thoại của tài khoản, yêu cầu đúng định dạng số điện thoại 10 chữ số, bắt đầu bằng số 09 / 03 / 05 / 07 / 08.
        \item Người dùng chọn để hệ thống gửi xác nhận tin nhắn OTP và chuyển tới trang xác nhận OTP.
    \end{enumerate}
\end{quote}
\begin{figure}[!htp]
    \centering
    \includegraphics[width=5in]{img/UI/customer/confirm_otp.png}
    \label{5}
    \newline
    \caption{Giao diện tạo lại mật khẩu mới}
\end{figure}
\textbf{Mô tả:}
\begin{quote}
    \begin{enumerate}
        \item Người dùng nhập mật khẩu mới cho tài khoản, yêu cầu tối thiểu 6 ký tự.
        \item Người dùng nhập lại mật khẩu mới cho tài khoản, yêu cầu phải giống với mật khẩu mới.
        \item Người dùng chọn "Xác nhận" để cập nhật lại mật khẩu mới cho tài khoản.
    \end{enumerate}
\end{quote}

\subsubsection{Header}
\begin{figure}[!htp]
    \centering
    \includegraphics[width=5in]{img/UI/customer/header.png}
    \label{6}
    \newline
    \caption{Giao diện phần header}
\end{figure}
\textbf{Mô tả:}
\begin{quote}
    \begin{enumerate}
        \item Chọn để chuyển tới trang chủ.
        \item Nhập tên sản phẩm cần tìm và nhấn enter để tìm kiếm sản phẩm.
        \item Chọn để chuyển đến trang quản lý lịch sử đơn hàng đã thực hiện.
        \item Chọn để chuyển đến trang quản lý giỏ hàng.
        \item Chọn để đăng nhập.
        \item Chọn để chuyển đến trang giới thiệu về hệ thống cửa hàng.
        \item Chọn để chuyển đến trang tất cả sản phẩm.
        \item Chọn để chuyển đến trang tất cả sản phẩm với lựa chọn lọc sẵn là sản phẩm dành cho nam.
        \item Chọn để chuyển đến trang tất cả sản phẩm với lựa chọn lọc sẵn là sản phẩm dành cho nữ.
        \item Chọn để chuyển đến trang tất cả sản phẩm với lựa chọn lọc sẵn là sản phẩm dành cho trẻ em.
        \item Chọn để chuyển đến trang tất cả sản phẩm với lựa chọn lọc sẵn là sản phẩm thuộc dòng phụ kiện.
        \item Chọn để chuyển đến trang các sản phẩm mới về.
        \item Chọn để chuyển đến trang các sản phẩm đang được flash sale.
        \item Chọn để chuyển đến trang hỗ trợ cho khách hàng.
    \end{enumerate}
\end{quote}

\subsubsection{Footer}
\begin{figure}[!htp]
    \centering
    \includegraphics[width=5in]{img/UI/customer/footer.png}
    \label{7}
    \newline
    \caption{Giao diện phần footer}
\end{figure}
\textbf{Mô tả:}
\begin{quote}
    \begin{enumerate}
        \item Chọn để chuyển tới trang chủ.
        \item Chọn để chuyển đến trang tất cả sản phẩm.
        \item Chọn để chuyển đến trang hỗ trợ cho khách hàng.
        \item Chọn để chuyển đến trang giới thiệu về hệ thống cửa hàng.
        \item Chọn để xem hướng dẫn chọn size.
        \item Chọn để chuyển đến trang quản lý lịch sử đơn hàng đã thực hiện.
    \end{enumerate}
\end{quote}

\newpage
\subsubsection{Trang chủ}
\begin{figure}[!htp]
    \centering
    \includegraphics[width=3.5in]{img/UI/customer/home.png}
    \label{8}
    \newline
    \caption{Giao diện trang chủ khi người dùng truy cập vào trang web}
\end{figure}
\textbf{Mô tả:}
\begin{quote}
    \begin{enumerate}
        \item Chọn để xem thông tin chi tiết của sản phẩm.
        \item Chọn để thêm nhanh sản phẩm vào giỏ hàng.
    \end{enumerate}
\end{quote}

\subsubsection{Thông tin chi tiết sản phẩm}
\begin{figure}[!htp]
    \centering
    \includegraphics[width=5in]{img/UI/customer/product_detail.png}
    \label{9}
    \newline
    \caption{Giao diện thông tin chi tiết của sản phẩm}
\end{figure}
\textbf{Mô tả:}
\begin{quote}
    \begin{enumerate}
        \item Chọn màu của sản phẩm mà người dùng muốn.
        \item Chọn size của sản phẩm mà người dùng muốn.
        \item Chọn để giảm số lượng mà người dùng muốn.
        \item Chọn để tăng số lượng mà người dùng muốn.
        \item Nhập để thay đổi số lượng mà người dùng muốn.
        \item Chọn để xem hướng dẫn chọn size cho sản phẩm.
        \item Chọn để thêm sản phẩm với màu sắc, size và số lượng mà người dùng đã chọn.
        \item Chọn để xem ảnh chi tiết mà người dùng muốn.
    \end{enumerate}
\end{quote}

\begin{figure}[!htp]
    \centering
    \includegraphics[width=5in]{img/UI/customer/guide_size.png}
    \label{10}
    \newline
    \caption{Giao diện hướng dẫn chọn size cho sản phẩm}
\end{figure}
\newpage

\subsubsection{Giỏ hàng}
\begin{figure}[!htp]
    \centering
    \includegraphics[width=5in]{img/UI/customer/cart.png}
    \label{11}
    \newline
    \caption{Giao diện quản lý giỏ hàng.}
\end{figure}
\textbf{Mô tả:}
\begin{quote}
    \begin{enumerate}
        \item Chọn để loại sản phẩm khỏi danh sách muốn đặt hàng.
        \item Chọn để xóa sản phẩm khỏi giỏ hàng.
        \item Chọn để giảm số lượng mà người dùng muốn.
        \item Chọn để tăng số lượng mà người dùng muốn.
        \item Nhập để thay đổi số lượng mà người dùng muốn.
        \item Chọn để thêm sản phẩm vào danh sách muốn đặt hàng.
        \item Nhập mã giảm giá.
        \item Chọn để áp dụng mã giảm giá.
        \item Chọn để thực hiện thanh toán để đặt hàng.
    \end{enumerate}
\end{quote}
\newpage



\subsubsection{Thanh toán}
\begin{figure}[!htp]
    \centering
    \includegraphics[width=5in]{img/UI/customer/payment.png}
    \label{12}
    \newline
    \caption{Giao diện thanh toán.}
\end{figure}
\textbf{Mô tả:}
\begin{quote}
    \begin{enumerate}
        \item Nhập tên người nhận hàng.
        \item Nhập số điện thoại người nhận hàng.
        \item Nhập địa chỉ email người nhận hàng.
        \item Nhập địa chỉ tỉnh/thành nhận hàng.
        \item Nhập địa chỉ quận/huyện nhận hàng.
        \item Nhập địa chỉ phường/xã nhận hàng.
        \item Nhập địa chỉ số nhà, tên đường nhận hàng.
        \item Chọn để thực hiện thanh toán trực tuyến.
        \item Chọn để thực hiện thanh toán trực tiếp khi nhận hàng.
        \item Chọn để thực hiện thanh toán trực tuyến thông qua VNPay.
        \item Chọn để thực hiện thanh toán trực tuyến thông qua Momo.
        \item Nhập để thêm ghi chú cho đơn hàng.
        \item Chọn để xác nhận đồng ý điều khoản mua hàng.
        \item Nhập mã giảm giá.
        \item Chọn để kiểm tra mã giảm giá và áp dụng giảm giá vào đơn hàng.
        \item Chọn để xác nhận đặt hàng, nếu chọn thanh toán trực tuyến thì sẽ được chuyển đến trang thanh toán của bên thứ ba mà người dùng chọn, nếu chọn thanh toán khi nhận hàng thì đơn hàng sẽ được hoàn tất khi tất cả thông tin hợp lệ.
    \end{enumerate}
\end{quote}

\subsubsection{Giao diện quản lý đơn hàng}
\begin{figure}[!htp]
    \centering
    \includegraphics[width=5in]{img/UI/customer/customer_order.png}
    \label{13}
    \newline
    \caption{Giao diện quản lý đơn hàng của khách hàng.}
\end{figure}
\textbf{Mô tả:}
\begin{quote}
    \begin{enumerate}
        \item Chọn để xem danh sách các đơn hàng đã nhận.
        \item Chọn để xem thông tin chi tiết của đơn hàng.
        \item Chọn để thực hiện thêm lại các sản phẩm của đơn hàng vào giỏ hàng.
    \end{enumerate}
\end{quote}

\newpage


\subsubsection{Giao diện chi tiết đơn hàng}
\begin{figure}[!htp]
    \centering
    \includegraphics[width=5in]{img/UI/customer/order_detail.png}
    \label{14}
    \newline
    \caption{Giao diện xem chi tiết đơn hàng của khách hàng.}
\end{figure}
\textbf{Mô tả:}
\begin{quote}
    \begin{enumerate}
        \item Chọn để xuất hóa đơn cho đơn hàng.
    \end{enumerate}
\end{quote}

\newpage

\subsubsection{Giao diện thông tin khách hàng}
\begin{figure}[!htp]
    \centering
    \includegraphics[width=5in]{img/UI/customer/customer_order.png}
    \label{15}
    \newline
    \caption{Giao diện xem thông tin tài khoản của khách hàng.}
\end{figure}
\textbf{Mô tả:}
\begin{quote}
    \begin{enumerate}
        \item Nhập để thay đổi tên của khách hàng.
        \item Nhập để thay đổi email của khách hàng.
        \item Chọn để thực hiện thay đổi mật khẩu của khách hàng.
        \item Nhập để thay đổi số nhà, đường của địa chỉ mặc định mà khách hàng muốn nhận hàng.
        \item Chọn để thay đổi phường/xã của địa chỉ mặc định mà khách hàng muốn nhận hàng.
        \item Chọn để thay đổi quận/huyện của địa chỉ mặc định mà khách hàng muốn nhận hàng.
        \item Chọn để thay đổi tỉnh/thành phố của địa chỉ mặc định mà khách hàng muốn nhận hàng.
        \item Chọn để thay đổi ảnh đại hình cho tài khoản của khách hàng.
    \end{enumerate}
\end{quote}

\newpage

Giao diện sau khi chọn "Đổi mật khẩu":
\begin{figure}[!htp]
    \centering
    \includegraphics[width=5in]{img/UI/customer/customer_order.png}
    \label{16}
    \newline
    \caption{Giao diện đổi mật khẩu cho tài khoản.}
\end{figure}

\textbf{Mô tả:}
\begin{quote}
    \begin{enumerate}
        \item Nhập để điền mật khẩu cũ, yêu cầu tối thiểu 6 ký tự.
        \item Nhập để điền mật khẩu mới, yêu cầu tối thiểu 6 ký tự.
        \item Nhập để điền xác nhận mật khẩu mới, yêu cầu giống mật khẩu mới.
        \item Chọn để xác nhận thay đổi mật khẩu.
        \item Chọn để tắt màn hình khi không muốn thay đổi mật khẩu nữa.
    \end{enumerate}
\end{quote}

\newpage



%%%%%%%%%%%%%%%%%%%%%%%%%%%%%%%%%%%%

\chapter{Công nghệ}
% technology
% % !TeX root = ..\main.tex

\section{Front-end}
\subsection{Next.js}

\hspace*{0.5cm} Next.js là một framework được phát triển bổ sung các khả năng render phía máy chủ (SSR). Next.js xây dựng dựa trên nền tảng là thư viện React, do đó các ứng dụng Next.js tương tự như các ứng dụng từ React nhưng được thêm các tính năng bổ sung.\\

\begin{figure}[!htp]
    \begin{center}
        \includegraphics[width=5cm]{img/Technology/next.png}
    \end{center}
    \caption{Logo Next.js \cite{technologyNextjs}}
\end{figure}

Ưu điểm của Next.js \cite{technologyNextAdvance}:
\begin{itemize}
    \item Dễ hiểu và dễ sử dụng: Có nền tảng từ ReactJs nên không khó để học và sử dụng với cộng đồng sử dụng lớn
    \item Hỗ trợ SEO: Next.js được bổ sung khả năng SSR, giúp hỗ trợ cho SEO tốt hơn nhiều so với React
\end{itemize}

Nhược điểm của Next.js \cite{technologyNextAdvance}::
\begin{itemize}
    \item Quá phức tạp và không cần thiết đối với các ứng dụng đơn giản
    \item Thời gian build hệ thống lớn hơn so với React
\end{itemize}

\subsection{Ant Design}
\hspace*{0.5cm} Ant Design là một thư viện của React cung cấp rất nhiều component có tính thông dụng cao, đa dạng và đầy đủ để đáp ứng các nhu cầu để không cần phải định nghĩa thêm các base component nữa. Ant Design có một cộng đồng sử dụng rất lớn, hiện tại số lượt started ở trên github của Ant Design đã đạt hơn 83 nghìn lượt \cite{technologyAntdStar}.

\begin{figure}[!htp]
    \begin{center}
        \includegraphics[width=5cm]{img/Technology/antd.jpg}
    \end{center}
    \caption{Logo Ant Design \cite{technologyAntd}}
\end{figure}

\hspace*{0.5cm} Ưu điểm của Ant Design bao gồm:
\begin{itemize}
    \item Số lượng base component: Như đã giới thiệu, Ant Design cung cấp rất nhiều base component đa dạng và đầy đủ giúp ta không cần xài thêm thư viện giao diện khác nữa.
    \item Khả năng customize base component: Ant Design cung cấp rất nhiều props để lập trình viên có thể tùy biến linh hoạt theo nhu cầu.
    \item Dễ học, dễ hiểu: Ant Design cung cấp bộ tài liệu hướng dẫn rất chi tiết cùng với demo tương ứng giúp người dùng có thể dễ dàng hiểu rõ và sử dụng.
    \item Responsive: Ant Design hỗ trợ dàn layout tương thích trên nhiều thiết bị và loại màn hình khác nhau, vì vậy không cần phải xài thêm các thư viện khác để responsive.
    \item Cộng đồng sử dụng: Ant Design có cộng đồng sử dụng lớn giúp cho việc tìm kiếm giải đáp thắc mắc rất dễ dàng.
\end{itemize}


\section{Back-end}
\hspace{0.5cm}Ở phần back-end, nhóm sử dụng ngôn ngữ Golang để hiện thực. Golang là một ngôn ngữ mới, được phát triển bởi Google nhằm sử dụng cho các hệ thống scale lớn. Đây là ngôn ngữ hiệu quả cho việc xây dựng hệ thống nói chung, và xây dựng hệ thống microservice nói riêng.

\begin{figure}[!htp]
    \begin{center}
        \includegraphics[width=5cm]{img/Technology/Golang.png}
    \end{center}
    \caption{Logo Golang \cite{technologyGolang}}
\end{figure}

Những lợi ích mà Golang có bao gồm\cite{technologyGolangReasons}:
\begin{itemize}
    \item Golang là ngôn ngữ static-typed, kiểu dữ liệu của biến là cố định và được khai báo từ đầu. Do đó các lỗi xảy ra liên quan đến kiểu dữ liệu đều được phát hiện và ngăn chặn ở bước compile. Ngoài ra tính năng thu gom rác (Garbage collection) giúp Golang quản lý bộ nhớ an toàn và hạn chế lỗi xuất hiện.
    \item Golang là ngôn ngữ đơn giản và dễ học. Cú pháp đơn giản và trực quan của Golang giúp cho ngôn ngữ này được ưa thích và người dùng có thể học dễ dàng.
    \item Golang là ngôn ngữ phổ biến, được sử dụng rộng rãi. Ngôn ngữ này được tạo nên và phát triển bởi Google, cùng cộng đồng lớn, do đó việc học và tìm sự trợ giúp khi sử dụng Golang trở nên dễ dàng hơn.
    \item Thời gian build và thực thi thấp hơn. Golang được tối ưu để thời gian build và chạy ít hơn.
    \item Golang hỗ trợ mạnh việc lập trình concurrency. Điều này giúp các ứng dụng Golang có khả năng mở rộng lớn.
\end{itemize}
\hspace{0.5cm}Qua đó, có thể thấy Golang có nhiều ưu điểm khi sử dụng. Với lợi thế về hiệu năng, khả năng mở rộng và độ tin cậy, Golang là lựa chọn thích hợp cho nhóm trong dự án này.

%%%%%%%%%%%%%%%%%%%%%%%%%%%%%%%%%%%%

\chapter{Tổng kết}
% summary
% % !TeX root = ..\main.tex
\section{Tổng kết}
%%%%%%%%%%%%%%%%%%%%%%%%%
\subsection{Kế hoạch cho giai đoạn luận văn}

\begin{figure}
    \begin{center}
        \includegraphics[width=14cm]{img/plan.png}
    \end{center}
    \caption{Biểu đồ Gantt cho kế hoạch đồ án tốt nghiệp}
\end{figure}

Lưu ý: Khoảng trống từ ngày 16/1 đến 29/1 là thời gian nghỉ tết

%%%%%%%%%%%%%%%%%%%%%%%%%
\subsection{Phân chia công việc}


%%%%%%%%%%%%%%%%%%%%%%%%%%%%%%%%%%%%

% refer
% !TeX root = ..\main.tex

\begin{thebibliography}{9}
    \bibitem{theoryBPMN0}
    Giải ngố các ký hiệu BPMN - thinhnotes.com - Truy cập lần cuối: 10/12/2022 - Tại: \url{https://thinhnotes.com/chuyen-nghe-ba/giai-ngo-cac-ky-hieu-bpmn/}

    \bibitem{theoryBPMN1}
    What is BPMN - Definition, Elements and Purpose(2022) - Allison Lynch - edrawsoft.com - Truy cập lần cuối: 10/12/2022 - Tại: \url{https://www.edrawsoft.com/what-is-bpmn.html?gclid=Cj0KCQiAvqGcBhCJARIsAFQ5ke5U9xsSqU24T_1MgVtsStQ3wV0NbkzXcIY-rDCU8JAet5rGQMc6PyIaAgdxEALw_wcB}

    \bibitem{theoryBPEL}
    A Hands-on Introduction to BPEL - Matjaz B. Juric - oracle.com - Truy cập lần cuối: 10/12/2022 - Tại: \url{https://www.oracle.com/technical-resources/articles/matjaz-bpel.html}

    \bibitem{theoryConvert}
    Translating BPMN to BPEL - Chun Ouyang \& Wil Van der Aalst \& Marlon Dumas \& Arthur Ter - Chương 3

    \bibitem{technologyFE}
    Tổng hợp top Front end Framework phổ biến nhất hiện nay (03/03/2022) - Niit Hanoi - niithanoi.edu.vn
    - Truy cập lần cuối: 02/05/2023 - Tại: \url{https://niithanoi.edu.vn/front-end-framework-pho-bien-nhat.html}

    \bibitem{technologyNextjs}
    NextJs Logo (12/12/2021) - Uttam Kumar Dash - ui-lib.com
    - Truy cập lần cuối: 02/05/2023 - Tại: \url{https://ui-lib.com/blog/nextjs-boilerplates/}


    \bibitem{technologyNextAdvance}
    React JS Web Development PROS AND CONS OF NEXTJS – 2022 UPDATED VERSION - By Tomasz Grabski - pagepro.co - Truy cập lần cuối: 01/05/2023 - Tại: \url{https://pagepro.co/blog/pros-and-cons-of-nextjs/}

    \bibitem{technologyAntd}
    Introduction - Ant Design - Ant Group and Ant Design Community - ant.design - Truy cập lần cuối: 10/12/2022 - Tại: \url{https://ant.design/docs/spec/introduce}

    \bibitem{technologyAntdStar}
    ant-design/ant-design - Github - github.com - Truy cập lần cuối: 10/12/2022 - Tại: \url{https://github.com/ant-design/ant-design}

    \bibitem{technologyAntdAdvance}
    Ant-design là cái quái gì? (19/08/2022) - @Cuongnp - viblo.asia - Truy cập lần cuối: 10/12/2022 - Tại: \url{https://viblo.asia/p/ant-design-la-cai-quai-gi-aWj53ewQZ6m}


    \bibitem{technologyGolang}
    Golang Logo (10/01/2020) - Lou Bichard - openupthecloud.com - Truy cập lần cuối: 15/12/2022 - Tại: \url{https://openupthecloud.com/cloud-engineers-code/golang/}

    \bibitem{technologyGolangReasons}
    Why You Should Go with Go for Your Next Software Project (09/2022) - Dimitar Grancharov - scalefocus.com - Truy cập lần cuối: 15/12/2022 - Tại: \url{https://www.scalefocus.com/blog/why-you-should-go-with-go-for-your-next-software-project}

    \bibitem{archChoreography}
    Fundamentals of Software Architecture An Engineering Approach - Mark Richards \& Neal Ford - Chương 17 - Trang 256

    \bibitem{archOrchestration}
    Fundamentals of Software Architecture An Engineering Approach - Mark Richards \& Neal Ford - Chương 17 - Trang 257



\end{thebibliography}

%%%%%%%%%%%%%%%%%%%%%%%%%%%%%%%%%%%%
\end{document}